% Αυτοσυνεπής Αναλογική Απόσβεση ανά Κόμβο για την Κατάταξη Προσβασιμότητας σε Στοχαστικά Γραφήματα
% Μεταγλώττιση:  xelatex paper_gr.tex  (τρεις φορές, έως ότου σταθεροποιηθούν οι παραπομπές)
% Κάθε πειραματική τιμή είναι μακροεντολή του generated/macros.tex ή πίνακας του
% generated/tables_gr.tex, που παράγονται από το ../make_paper_assets.py με βάση τα ../results/.
\documentclass[11pt,a4paper]{article}
\usepackage[a4paper,top=2.5cm,bottom=2.5cm,left=2.2cm,right=2.2cm]{geometry}
\usepackage{amsmath,amsthm}
\usepackage{mathptmx}
\usepackage{amssymb}
\usepackage[no-math]{fontspec}
\usepackage[english,main=greek]{babel}
\babelfont{rm}[Language=Default,RawFeature={-ccmp}]{FreeSerif}
\babelfont{sf}[Language=Default,RawFeature={-ccmp}]{FreeSans}
\babelfont{tt}[Language=Default,Scale=0.86]{DejaVu Sans Mono}
\usepackage{microtype}
\usepackage{booktabs}
\usepackage{array}
\usepackage{graphicx}
\graphicspath{{../figures_gr/}}
\usepackage[font=small,labelfont=bf,skip=4pt]{caption}
\usepackage{enumitem}
\setlist{itemsep=2pt,topsep=3pt}
\usepackage{cite}
\usepackage{xurl}
\usepackage[hidelinks]{hyperref}
\usepackage{placeins}

\emergencystretch=2em
\setlength{\parskip}{0pt}
\renewcommand{\floatpagefraction}{0.8}
\renewcommand{\topfraction}{0.9}
\renewcommand{\textfraction}{0.08}

\newtheorem{theorem}{Θεώρημα}
\newtheorem{lemma}{Λήμμα}
\newtheorem{corollary}{Πόρισμα}
\theoremstyle{definition}
\newtheorem{definition}{Ορισμός}
\newtheorem{remark}{Παρατήρηση}
\theoremstyle{plain}
\newcommand{\genpropname}{}
\newtheorem*{genprop}{\genpropname}
\newenvironment{prop}[1]{\renewcommand{\genpropname}{Πρόταση #1}\begin{genprop}}{\end{genprop}}

\newcommand{\Pesc}{P_{\mathrm{esc}}}
\newcommand{\Pescbar}{\bar{P}_{\mathrm{esc}}}
\newcommand{\clip}{\mathrm{clip}}
\newcommand{\diag}{\mathrm{diag}}
\newcommand{\E}{\mathbb{E}}
\newcommand{\N}{\mathcal{N}}
\newcommand{\AP}{\mathrm{AP}}
\newcommand{\one}{\mathbf{1}}

\input{generated/macros.tex}
\input{generated/tables_gr.tex}

\title{\textbf{Αυτοσυνεπής Αναλογική Απόσβεση ανά Κόμβο για την Κατάταξη Προσβασιμότητας σε Στοχαστικά Γραφήματα: Ακριβής Θεωρία, Δυναμικά Εμπόδια και Υβριδική Εξομάλυνση Διπλού Τελεστή}}
\author{Ιωάννης Κωνσταντέλλος\\[2pt]
\small Τμήμα Μηχανικών Ηλεκτρονικών Υπολογιστών και Πληροφορικής\\
\small Πανεπιστήμιο Πατρών}
\date{Οκτώβριος 2026}

\begin{document}
\maketitle

\begin{abstract}
\noindent Οι αποσβεννύμενες διαχύσεις κατατάσσουν τις κορυφές ενός στοχαστικού γραφήματος $G=(V,E,p)$ ως προς την αναμενόμενη προσβάσιμη μάζα $y^*_i=\E\big[\sum_{v\in C(i)}x_v\big]$ με κόστος λίγων γινομένων αραιού πίνακα επί διάνυσμα, εφαρμόζουν όμως τον ίδιο συντελεστή απόσβεσης σε όλες τις κορυφές. Η προηγούμενη εργασία μας βαθμονόμησε τον συντελεστή αυτόν από τη μέση πιθανότητα διαφυγής~\cite{konstantellos26} και ανέφερε ότι ένα διάνυσμα απόσβεσης ανά κορυφή, προσαρμοσμένο με βέλτιστες παραμέτρους αναφοράς (oracle / μαντείο), κερδίζει επιπλέον \PrevHeadAll{} στον συντελεστή $\rho$ του Spearman (\PrevHeadGrid{} στα πλέγματα). Η παρούσα εργασία μελετά έναν τρόπο καθορισμού της απόσβεσης ανά κορυφή χωρίς ελεύθερες παραμέτρους, την Αρχή της Αναλογικής Σχετικής Επιρροής $\alpha_i/(1-\alpha_i)=(y_i-x_i)/x_i$, συζευγμένη με το βήμα διάχυσης $y_i=(1-\alpha_i)x_i+\alpha_iS_i$. Το συζευγμένο σύστημα παραγοντοποιείται ακριβώς ως $(y_i-x_i)(y_i+x_i-S_i)=0$ (Θεώρημα~1). Αποδεικνύουμε έξι ιδιότητες των υλοποιήσεών του: κατάρρευση στην αρχική κατανομή σε πεπερασμένο χρόνο για υποστοχαστικούς τελεστές, απορρόφηση με τετραγωνική συστολή για τη χρονικά υστερημένη μορφή, αστάθεια κάθε θετικής αυτοσυνεπούς λύσης, συστολή μέσου Lehmer με σταθερό σημείο κάτω από την αρχική κατανομή, ένα εμπόδιο προσήμου εξαιτίας του οποίου ο νόμος υποβαθμίζει τις καλά συνδεδεμένες κορυφές όπου το σήμα των γειτόνων τους είναι μικρότερο από την αρχική τους μάζα, και εκφυλισμό της παραλλαγής έλξης υπό ομοιόμορφες αρχικές κατανομές· μια ακριβής ταυτότητα λογαριθμικής διασποράς αντικαθιστά τον ευρετικό ισχυρισμό περί ομοιομορφίας της αρχικής κατανομής. Μια προκαταχωρισμένη αξιολόγηση σε \NTotal{} στιγμιότυπα, με όλες τις διαμορφώσεις παγωμένες σε σύνολο ανάπτυξης και διόρθωση Holm επί \NFrozen{} διαμορφώσεων, επιβεβαιώνει τη θεωρία: έναντι του βαθμονομημένου βαθμωτού κανόνα, ο αγκυρωμένος νόμος μεταβάλλει το $\rho$ κατά \TestAPoneD, η υβριδική εκδοχή όπως διατυπώθηκε κατά \TestPBXnomBD{} και ο νόμος με το αληθές $y^*$ κατά \TestTruePropD. Εκτός δείγματος, το περιθώριο βελτίωσης της απόσβεσης ανά κόμβο είναι \HeadAll. Κέρδος αποφέρει η υβριδική αρχιτεκτονική όταν χρησιμοποιείται ως συνδυασμός διπλού τελεστή: ο συνδυασμός της βαθμονομημένης διάχυσης με λίγα βήματα διάδοσης στον ακατέργαστο πίνακα πιθανοτήτων γειτνίασης βελτιώνει το $\rho$ κατά \TestPBXbestD{} έως \TestLinNineD{} σε ανεξάρτητα στιγμιότυπα και μεταφέρεται σε πραγματικά δίκτυα και στο $n=900$. Αντιστοιχισμένοι έλεγχοι αποδίδουν το κέρδος στη χαμηλή συσχέτιση των σφαλμάτων τάξης των δύο τελεστών και όχι στον αναλογικό νόμο. Διορθώνουμε επίσης δύο ισχυρισμούς του~\cite{konstantellos26}: η βαθμονόμηση εξαρτάται από την αρχική κατανομή, και ένας δειγματολήπτης Monte Carlo κατά δέσμες είναι αρκετά γρήγορος ώστε ο ντετερμινιστικός συνδυασμός να είναι φθηνότερος από ισοδύναμης ακρίβειας δειγματοληψία μόνο πέρα από ένα μέγεθος γραφήματος που εξαρτάται από την οικογένεια.
\end{abstract}

\section{Εισαγωγή}\label{sec:intro}

Έστω $G=(V,E)$ ένα μη κατευθυνόμενο γράφημα με $n$ κορυφές, στο οποίο κάθε ακμή $(i,j)\in E$ υπάρχει ανεξάρτητα με πιθανότητα $p_{ij}\in(0,1]$, και έστω $x\in\mathbb{R}^n_{>0}$ με $\sum_ix_i=1$ μια αρχική κατανομή (prior) επί των κορυφών. Σε έναν δειγματοληπτημένο κόσμο $G'\sim G$, έστω $C(i)$ η συνεκτική συνιστώσα της κορυφής $i$. Η \emph{αναμενόμενη προσβάσιμη μάζα}
\begin{equation}\label{eq:target}
y^*_i=\E_{G'\sim G}\Big[\sum_{v\in C(i)}x_v\Big]
\end{equation}
είναι η εξάπλωση επιρροής από έναν σπόρο στο μοντέλο ανεξάρτητων καταρρακτών με βάρη κόμβων~\cite{kkt03} και, υπό ομοιόμορφη αρχική κατανομή, η κεντρικότητα συνιστώσας διήθησης ακμών (EPC)~\cite{hubba08,cytohubba14} διαιρεμένη με $n$. Ο ακριβής υπολογισμός της είναι \#P-δύσκολος, αφού περιέχει την αξιοπιστία δικτύου δύο τερματικών~\cite{provanball83}· στην πράξη εκτιμάται με δειγματοληψία κόσμων Monte Carlo, με μετάδοση μηνυμάτων~\cite{knz14} ή με φθηνούς ντετερμινιστικούς ευρετικούς αλγορίθμους που αρκεί να \emph{κατατάσσουν} σωστά τις κορυφές. Οι αποσβεννύμενες διαχύσεις της οικογένειας του εξατομικευμένου PageRank~\cite{pagerank99,jeh03,haveliwala03,gleich15} είναι οι συνηθέστεροι ευρετικοί αλγόριθμοι αυτού του τύπου.

Μια αποσβεννύμενη διάχυση αναμειγνύει την αρχική κατανομή με ένα διαδιδόμενο σήμα, $y\leftarrow(1-\alpha)x+\alpha Ty$, και ο συντελεστής απόσβεσης $\alpha$ καθορίζει πόσο μακριά ταξιδεύει η μάζα. Στο~\cite{konstantellos26} δείξαμε ότι το $\alpha$ δεν χρειάζεται να κληρονομείται από την αναζήτηση στον Ιστό ούτε να ρυθμίζεται έναντι της πραγματικής τιμής αναφοράς: ο αποκομμένος γραμμικός (clipped-linear) κανόνας
\begin{equation}\label{eq:rule}
\alpha_0=\clip\big(0.502+0.558\,\Pescbar,\;0.4,\;1\big),\qquad \Pesc(i)=1-\prod_{j\in N(i)}(1-p_{ij}),\qquad \Pescbar=\tfrac1n\textstyle\sum_i\Pesc(i),
\end{equation}
ανακτά το \PrevPctOracle\% της απόδοσης των βέλτιστων βαθμωτών παραμέτρων αναφοράς (oracle / μαντείο) ανά στιγμιότυπο. Η ίδια μελέτη ανέφερε ότι ένα διάνυσμα απόσβεσης \emph{ανά κορυφή}, προσαρμοσμένο με ανάβαση κατά συντεταγμένες, υπερέχει της καλύτερης βαθμωτής απόσβεσης κατά \PrevHeadAll{} στο $\rho$ του Spearman συνολικά και κατά \PrevHeadGrid{} στα κανονικά πλέγματα (grids), όπου η πιθανότητα διαφυγής ενός βήματος είναι ίδια για όλες τις εσωτερικές κορυφές. Αν έστω και μέρος αυτού του περιθωρίου μπορούσε να αξιοποιηθεί χωρίς επίβλεψη, μια διάχυση θα μπορούσε να προσαρμόζει την απόσβεσή της σε κάθε κορυφή.

Τα συνοδευτικά χειρόγραφα της παρούσας εργασίας πρότειναν έναν τέτοιο μηχανισμό, την \emph{Αρχή της Αναλογικής Σχετικής Επιρροής}: σε κάθε κορυφή, ο λόγος πιθανοτήτων διάδοσης έναντι διατήρησης της μάζας πρέπει να ισούται με την εξωτερική επιρροή της κορυφής σε σχέση με τη δική της αρχική μάζα,
\begin{equation}\label{eq:system}
\begin{cases}
y_i=(1-\alpha_i)\,x_i+\alpha_i\,S_i & \text{(εξίσωση διάχυσης)},\\[3pt]
\dfrac{\alpha_i}{1-\alpha_i}=\dfrac{y_i-x_i}{x_i} & \text{(εξίσωση αναλογίας)},
\end{cases}
\end{equation}
όπου $S_i=(Ty)_i$ είναι το σήμα των γειτόνων για έναν τελεστή διάδοσης $T$. Το σύστημα επιλύεται σε κλειστή μορφή, υποδεικνύει τρεις επαναληπτικές υλοποιήσεις και μπορεί να εκτελεστεί είτε από την αρχή (Εκδοχή~(α)) είτε ως στάδιο «εξομάλυνσης» (polishing, «λιμάρισμα») μετά από μια βαθμονομημένη διάχυση βάσης (Εκδοχή~(β)).

Η εργασία αυτή εξετάζει την αρχή όσο το δυνατόν πιο αυστηρά, μέσω ακριβούς ανάλυσης και μέσω ενός προκαταχωρισμένου πειράματος, και αναφέρει τι επιβιώνει. Το αποτέλεσμα είναι σε μεγάλο βαθμό αρνητικό, και το αρνητικό σκέλος εξηγείται: ο στόχος~\eqref{eq:target} δεν ικανοποιεί το συζευγμένο σύστημα σε κανέναν κανονικοποιημένο τελεστή, οι υλοποιήσεις του είτε καταρρέουν στην αρχική κατανομή είτε απομακρύνονται από το αυτοσυνεπές σημείο τους, και ο νόμος ανεβάζει ή κατεβάζει μια κορυφή ανάλογα με το πρόσημο του $S_i-x_i$ και όχι ανάλογα με την προσβασιμότητά της. Το πείραμα επιβεβαιώνει καθέναν από αυτούς τους μηχανισμούς. Η αρχιτεκτονική της Εκδοχής~(β), ωστόσο, αποδεικνύεται χρήσιμη για έναν λόγο που η αρχή δεν είχε προβλέψει: τον συνδυασμό δύο τελεστών διάδοσης των οποίων τα σφάλματα είναι ασθενώς συσχετισμένα. Οι συνεισφορές μας είναι οι εξής:
\begin{enumerate}[label=(\roman*)]
\item \textbf{Ακριβής άλγεβρα και μια εμφωλευμένη γενίκευση} (Ενότητα~\ref{sec:principle}). Αποδεικνύουμε την παραγοντοποίηση του συζευγμένου συστήματος (Θεώρημα~\ref{thm:factor}), προσδιορίζουμε τι απαιτεί από την κλίμακα του τελεστή, ορίζουμε με ακρίβεια τις υλοποιήσεις και εντάσσουμε τον νόμο σε μια αγκυρωμένη οικογένεια δυνάμεων $\AP(\eta)$ που περιέχει τον βαθμονομημένο βαθμωτό κανόνα για $\eta=0$.
\item \textbf{Δυναμικά εμπόδια} (Ενότητα~\ref{sec:theory}). Οι Προτάσεις~A--F θεμελιώνουν κατάρρευση σε πεπερασμένο χρόνο, απορρόφηση, αστάθεια με φασματική ακτίνα σε κλειστή μορφή, συστολή μέσου Lehmer κάτω από την αρχική κατανομή, ένα εμπόδιο προσήμου (με αυστηρή εκδοχή για γραμμικές διαχύσεις ανά κόμβο) και τον εκφυλισμό της παραλλαγής έλξης. Η Πρόταση~$2'$ μετατρέπει την υπόθεση της ομοιομορφίας της αρχικής κατανομής σε ακριβή ταυτότητα.
\item \textbf{Μια προκαταχωρισμένη αξιολόγηση} (Ενότητες~\ref{sec:protocol}--\ref{sec:results}) σε \NTotal{} στιγμιότυπα πέντε συνόλων, με \CfgTotal{} διαμορφώσεις που αναζητήθηκαν σε σύνολο ανάπτυξης, \NFrozen{} παγωμένες πριν από την εξέταση οποιουδήποτε ανεξάρτητου αποτελέσματος, ελέγχους κατά ζεύγη με διόρθωση Holm και αντιστοιχισμένους ελέγχους που απομονώνουν τη συνεισφορά του νόμου ανά κόμβο.
\item \textbf{Μια νέα μέτρηση του περιθωρίου βελτίωσης ανά κόμβο}: \HeadAll{} εκτός δείγματος (\HeadApparent{} εντός δείγματος), \HeadGrid{} στα πλέγματα, συγκεντρωμένο στο καθεστώς πυκνών γραφημάτων με ετερογενή αρχική κατανομή.
\item \textbf{Εξομάλυνση διπλού τελεστή}: ένα εποικοδομητικό αποτέλεσμα της Εκδοχής~(β) που ανακτά το \TestLinNinePctOracle\% του βαθμωτού μαντείου σε ανεξάρτητα στιγμιότυπα, μαζί με τις ενδείξεις ότι το κέρδος της προέρχεται από τη συμπληρωματικότητα των τελεστών.
\item \textbf{Δύο διορθώσεις του~\cite{konstantellos26}} --- η βαθμονόμηση εξαρτάται από την αρχική κατανομή, και η δειγματοληψία Monte Carlo κατά δέσμες είναι πολύ φθηνότερη από ό,τι είχε υποτεθεί --- και ένα μετρημένο, εξαρτώμενο από την οικογένεια σημείο διασταύρωσης κόστους μεταξύ ντετερμινιστικής κατάταξης και δειγματοληψίας.
\end{enumerate}

Η Ενότητα~\ref{sec:related} επισκοπεί τη σχετική βιβλιογραφία, η Ενότητα~\ref{sec:principle} αναπτύσσει την αρχή και την άλγεβρά της, η Ενότητα~\ref{sec:theory} τη δυναμική θεωρία, η Ενότητα~\ref{sec:protocol} το πειραματικό πρωτόκολλο και η Ενότητα~\ref{sec:results} τα αποτελέσματα. Η Ενότητα~\ref{sec:discussion} συζητά τις συνέπειές τους, η Ενότητα~\ref{sec:limitations} τις απειλές για την εγκυρότητα και η Ενότητα~\ref{sec:conclusion} συμπεραίνει.

\section{Σχετική βιβλιογραφία}\label{sec:related}

\paragraph{Προσβασιμότητα σε στοχαστικά γραφήματα.} Τα στοχαστικά (πιθανοτικά) γραφήματα με ανεξάρτητες ακμές αποτελούν το καθιερωμένο μοντέλο για αναξιόπιστους συνδέσμους και αβέβαιες αλληλεπιδράσεις· ερωτήματα όπως οι πλησιέστεροι γείτονες~\cite{potamias10} και η συνεκτικότητα απαιτούν συλλογισμό επί εκθετικά πολλών κόσμων, και η αξιοπιστία δύο τερματικών είναι \#P-πλήρης~\cite{provanball83}. Η αναμενόμενη προσβάσιμη μάζα ταυτίζεται με την εξάπλωση επιρροής από έναν σπόρο στο μοντέλο ανεξάρτητων καταρρακτών~\cite{kkt03}, για την οποία οι εκτιμητές δειγματοληψίας κυριαρχούν στη μεγιστοποίηση επιρροής~\cite{borgs14,tim14}. Η μετάδοση μηνυμάτων για τη διήθηση σε αραιά δίκτυα~\cite{knz14,newman23} δίνει ακριβή στατιστικά συνιστωσών σε τοπικά δενδροειδή γραφήματα, και η πιθανότητα διαφυγής $\Pesc$ που χρησιμοποιήθηκε στο~\cite{konstantellos26} είναι ακριβώς η μηδενική επανάληψή της χωρίς συσχετίσεις. Η κεντρικότητα συνιστώσας διήθησης ακμών~\cite{hubba08,cytohubba14} είναι η ειδική περίπτωση του~\eqref{eq:target} για ομοιόμορφη αρχική κατανομή.

\paragraph{Ο συντελεστής απόσβεσης.} Το PageRank~\cite{pagerank99}, το εξατομικευμένο και θεματικά ευαίσθητο PageRank~\cite{jeh03,haveliwala03} και οι πολυάριθμες παραλλαγές τους~\cite{gleich15} μοιράζονται την παράμετρο απόσβεσης (ή επανεκκίνησης) $\alpha$. Η επίδρασή της στις κατατάξεις έχει μελετηθεί ως συνάρτηση~\cite{boldi05} και ως πηγή αστάθειας της κατάταξης~\cite{bressan10}· οι Constantine και Gleich~\cite{constantine10} θεωρούν το $\alpha$ τυχαία μεταβλητή και λαμβάνουν τον μέσο όρο ως προς αυτήν, ενώ αναλύσεις μέσου πεδίου περιγράφουν το εξατομικευμένο PageRank σε μοντέλα τυχαίων γραφημάτων~\cite{avrachenkov18}. Ο προηγούμενος κανόνας βαθμονόμησής μας~\cite{konstantellos26} προβλέπει το καλύτερο βαθμωτό $\alpha$ για την κατάταξη προσβασιμότητας από το $\Pescbar$.

\paragraph{Απόσβεση και επανεκκίνηση εξαρτώμενες από τον κόμβο.} Η εξάρτηση της πιθανότητας επανεκκίνησης από την κορυφή έχει μελετηθεί από διάφορες οπτικές. Οι Avrachenkov, van der Hofstad και Sokol~\cite{avrachenkov14} εισήγαγαν δύο γενικεύσεις του εξατομικευμένου PageRank με επανεκκίνηση εξαρτώμενη από τον κόμβο: η μία βασίζεται στο ποσοστό επισκέψεων σε κάθε κορυφή πριν από την επανεκκίνηση και η άλλη στην κορυφή που επισκέφθηκε ο περίπατος αμέσως πριν από αυτήν· οι δύο ταυτίζονται για σταθερή πιθανότητα επανεκκίνησης, και αμφότερες συνδέουν το ευθύ με το αντίστροφο εξατομικευμένο PageRank. Οι μερικώς απορροφητικοί τυχαίοι περίπατοι~\cite{wu12} αποδίδουν σε κάθε κορυφή τον δικό της ρυθμό απορρόφησης. Πιθανότητες επανεκκίνησης έχουν μαθευτεί με επίβλεψη για κατάταξη και πρόβλεψη συνδέσμων~\cite{jung19}, ενώ οι μέθοδοι γενικευμένου PageRank μαθαίνουν συντελεστές διάχυσης ανά βήμα για επέκταση σπόρων και για νευρωνικά δίκτυα γραφημάτων~\cite{li19gpr,gprgnn21,berberidis19}· ο δείκτης κύρους του Katz~\cite{katz53} είναι ο πρόγονος χωρίς απόσβεση αυτών των περικομμένων περιπάτων. Σε όλες αυτές τις εργασίες οι παράμετροι ανά κορυφή (ή ανά βήμα) είναι εξωγενείς: καθορίζονται σχεδιαστικά ή μαθαίνονται με επίβλεψη. Η αρχή που μελετάμε εδώ καθιστά το $\alpha_i$ \emph{ενδογενή} συνάρτηση της ίδιας της βαθμολογίας που υπολογίζεται. Αυτή η ανάδραση παράγει τόσο την ακριβή άλγεβρα της Ενότητας~\ref{sec:principle} όσο και τα δυναμικά εμπόδια της Ενότητας~\ref{sec:theory}.

\paragraph{Συνδυασμός κατατάξεων.} Ο μέσος όρος βαθμολογιών και τάξεων είναι οι απλούστερες μέθοδοι συνάθροισης κατατάξεων~\cite{dwork01}· οι συνδυασμοί μοντέλων ωφελούν όταν τα σφάλματα των μελών τους είναι ασθενώς συσχετισμένα~\cite{krogh95}. Η εξομάλυνση διπλού τελεστή που προτείνουμε είναι ένας συνδυασμός δύο μελών αυτού του είδους, και μετράμε απευθείας τη συσχέτιση των σφαλμάτων των μελών του.

\section{Η αρχή και η ακριβής άλγεβρά της}\label{sec:principle}

\subsection{Τελεστές και βαθμονομημένη διάχυση}\label{sec:operators}

Συμβολίζουμε με $N(i)$ τους γείτονες της $i$ και με $\Pesc(i)$ την πιθανότητα διαφυγής της~\eqref{eq:rule}, δηλαδή την πιθανότητα να υπάρχει τουλάχιστον μία από τις ακμές της. Ακολουθώντας το~\cite{konstantellos26}, χρησιμοποιούμε τον στοχαστικό κατά γραμμές πίνακα $W$ με $w_{ij}\propto-\ln(1-p_{ij})$ στις ακμές ($\sum_jw_{ij}=1$ για μη απομονωμένη $i$), τον τελεστή έλξης (pull) με κλιμάκωση διαφυγής $EW=\diag(\Pesc)\,W$, τον τελεστή ώθησης (push) $W^{\top}$ του κλασικού εξατομικευμένου PageRank, τον ακατέργαστο πίνακα πιθανοτήτων γειτνίασης $P=(p_{ij})$ και την εκδοχή του με κλιμάκωση διαφυγής $EP=\diag(\Pesc)\,P$. Οι γραμμές του $W$ αθροίζουν στη μονάδα και του $EW$ στο $\Pesc(i)\le1$· οι γραμμές του $P$ αθροίζουν στον αναμενόμενο βαθμό $\sum_jp_{ij}$, ο οποίος μπορεί να υπερβαίνει τη μονάδα.

Για ένα διάνυσμα απόσβεσης $\alpha\in[0,1]^n$ και $D=\diag(\alpha)$, η \emph{κανονικοποιημένη αποσβεννύμενη διάχυση} είναι το διάνυσμα Perron του
\begin{equation}\label{eq:diffusion}
B(\alpha)=(I-D)\,x\,\one^{\top}+D\,T,
\end{equation}
δηλαδή το σταθερό σημείο της $y\leftarrow\N\big((1-\alpha)\circ x+\alpha\circ Ty\big)$ με $\N(v)=v/\one^{\top}v$. Το υπολογίζουμε με την οκνηρή (lazy) μέθοδο δυνάμεων $y\leftarrow\N\big(\tfrac12(y+B(\alpha)y)\big)$, η οποία έχει το ίδιο σταθερό σημείο και συγκλίνει και σε διμερή γραφήματα. Με $T=EW$ και το βαθμωτό $\alpha_0$ της~\eqref{eq:rule}, αυτός είναι ο \emph{βαθμονομημένος κανόνας} του~\cite{konstantellos26}, η μέθοδος αναφοράς της παρούσας εργασίας. Ένα διάνυσμα βαθμολογιών αξιολογείται με τον συντελεστή $\rho$ του Spearman ως προς το~\eqref{eq:target}· οι ισόπαλες βαθμολογίες λαμβάνουν μέσες τάξεις και ένα σταθερό διάνυσμα βαθμολογιών λαμβάνει $\rho=0$.

\subsection{Η αρχή και το Θεώρημα~1}\label{sec:theorem}

\begin{theorem}[ακριβής λύση του συζευγμένου συστήματος]\label{thm:factor}
Έστω $x_i>0$ και $S_i\ge0$. Ένα ζεύγος $(y_i,\alpha_i)$ με $y_i>0$ και $\alpha_i\neq1$ λύνει το~\eqref{eq:system} αν και μόνο αν $\alpha_i=1-x_i/y_i$ και
\begin{equation}\label{eq:factor}
(y_i-x_i)\,(y_i+x_i-S_i)=0 .
\end{equation}
Οι λύσεις είναι ο τετριμμένος κλάδος $(y_i,\alpha_i)=(x_i,0)$ και, όταν $S_i>x_i$, ο ενεργός κλάδος
\begin{equation}\label{eq:active}
y_i=S_i-x_i,\qquad \alpha_i=\frac{S_i-2x_i}{S_i-x_i}=1-\frac{x_i}{S_i-x_i}.
\end{equation}
Ο ενεργός κλάδος έχει $\alpha_i\in(0,1)$ αν και μόνο αν $S_i>2x_i$· οι δύο κλάδοι συμπίπτουν για $S_i=2x_i$, όπου μηδενίζεται η διακρίνουσα $(S_i-2x_i)^2$.
\end{theorem}

\begin{proof}
Προσθέτοντας τη μονάδα και στα δύο μέλη της εξίσωσης αναλογίας λαμβάνουμε $1/(1-\alpha_i)=y_i/x_i$, άρα $\alpha_i=1-x_i/y_i$ (εδώ χρησιμοποιούνται τα $\alpha_i\ne1$ και $y_i\neq0$), και αντίστροφα η τιμή αυτή του $\alpha_i$ ικανοποιεί την εξίσωση αναλογίας. Αντικαθιστώντας την στην εξίσωση διάχυσης και πολλαπλασιάζοντας με $y_i>0$ λαμβάνουμε $y_i^2=x_i^2+(y_i-x_i)S_i$, δηλαδή $(y_i-x_i)(y_i+x_i)-S_i(y_i-x_i)=0$, που είναι η~\eqref{eq:factor}. Η ρίζα $y_i=x_i$ δίνει $\alpha_i=0$. Η ρίζα $y_i=S_i-x_i$ είναι θετική ακριβώς όταν $S_i>x_i$ και δίνει $\alpha_i=1-x_i/(S_i-x_i)=(S_i-2x_i)/(S_i-x_i)$, που ανήκει στο $(0,1)$ ακριβώς όταν $S_i>2x_i$. Το ανάπτυγμα της~\eqref{eq:factor} είναι $y_i^2-S_iy_i+x_iS_i-x_i^2=0$, με διακρίνουσα $S_i^2-4x_iS_i+4x_i^2=(S_i-2x_i)^2$.
\end{proof}

Το Θεώρημα~\ref{thm:factor} αποτελεί τον αλγεβρικό πυρήνα των συνοδευτικών χειρογράφων και είναι ορθό: στην υλοποίησή μας και οι δύο κλάδοι ικανοποιούν και τις δύο εξισώσεις με σχετικό υπόλοιπο \ThmResidual. Η ερμηνεία του, ωστόσο, χρειάζεται προσοχή.

\begin{remark}[τι απαιτεί το σύστημα από τον τελεστή]\label{rem:scale}
Στον ενεργό κλάδο η βαθμολογία $y_i=S_i-x_i$ \emph{μειώνεται} ως προς την αρχική μάζα της ίδιας της κορυφής για σταθερό σήμα γειτόνων, ενώ ο στόχος $y^*_i=x_i+\E\big[\sum_{v\in C(i)\setminus\{i\}}x_v\big]$ αυξάνεται ως προς αυτήν. Ο στόχος ικανοποιεί το συζευγμένο σύστημα στον ενεργό κλάδο μόνο αν $Ty^*=y^*+x$, δηλαδή μόνο αν ο τελεστής απεικονίζει τον στόχο στον εαυτό του συν την αρχική κατανομή. Αφού $y^*+x>y^*$ κατά συνιστώσα, το φράγμα Collatz--Wielandt~\cite{hornjohnson13} επιβάλλει τότε $\rho(T)>1$: κανένας υποστοχαστικός τελεστής, και ειδικότερα καμία από τις κανονικοποιημένες διαχύσεις που χρησιμοποιούνται στην πράξη, δεν μπορεί να ικανοποιεί το σύστημα στον στόχο. Το συζευγμένο σύστημα είναι επομένως περιορισμός επί της κλίμακας του τελεστή και όχι ιδιότητα της προσβασιμότητας.
\end{remark}

\subsection{Τρεις υλοποιήσεις και οι κλίμακές τους}\label{sec:realizations}

Ένας επαναληπτικός επιλυτής υπολογίζει το $S^{(t)}=Ty^{(t)}$ και ενημερώνει όλες τις κορυφές παράλληλα. Ακολουθώντας τα χειρόγραφα, μελετάμε τρεις υλοποιήσεις, όλες της μορφής $y^{(t+1)}=R\big((1-\alpha^{(t+1)})\circ x+\alpha^{(t+1)}\circ S^{(t)}\big)$, για μια απεικόνιση κλίμακας $R$ που εξετάζεται παρακάτω:
\begin{description}[leftmargin=1.2em,style=nextline]
\item[Υλοποίηση Α (ταυτόχρονη ρίζα).] $\alpha_i^{(t+1)}=\clip\big((S_i^{(t)}-2x_i)/(S_i^{(t)}-x_i),0,\alpha_{\max}\big)$ αν $S_i^{(t)}>2x_i$ και $0$ διαφορετικά. Χωρίς το άνω όριο, πρόκειται για την άνω περιβάλλουσα $y_i^{(t+1)}=\max\big(x_i,\,S_i^{(t)}-x_i\big)$ του Θεωρήματος~\ref{thm:factor}.
\item[Υλοποίηση Β (χρονικά υστερημένη).] $\alpha_i^{(t+1)}=1-x_i/y_i^{(t)}$, υλοποιημένη ως $(y_i^{(t)}-x_i)_+/\big((y_i^{(t)}-x_i)_++x_i\big)$ ώστε το $\alpha$ να μην είναι ποτέ αρνητικό· επομένως $y_i^{(t+1)}=x_i^2/y_i^{(t)}+(1-x_i/y_i^{(t)})S_i^{(t)}$ όσο $y_i^{(t)}\ge x_i$. Τα σταθερά σημεία της είναι ακριβώς οι λύσεις της~\eqref{eq:factor}.
\item[Υλοποίηση Γ (άμεσο σήμα γειτόνων).] $\alpha_i^{(t+1)}=\gamma S_i^{(t)}/\big(x_i+\gamma S_i^{(t)}\big)$· για το $\gamma=1$ των χειρογράφων η ενημέρωση είναι ο μέσος Lehmer $y_i^{(t+1)}=L_2(x_i,S_i^{(t)})=(x_i^2+S_i^{(t)2})/(x_i+S_i^{(t)})$~\cite{lehmer71}.
\end{description}
Όλες οι υλοποιήσεις ξεκινούν από τη θερμή εκκίνηση ενός βήματος $y^{(0)}=x+Tx$. Επειδή ο νόμος συγκρίνει το $y_i$ με το $x_i$, η κλίμακα του $y$ έχει σημασία, και μελετάμε τέσσερις συμβάσεις για την $R$: την \emph{ακατέργαστη} κλίμακα ($R$ η ταυτοτική), την κλίμακα \emph{$L_1$} ($R=\N$) και την κλίμακα \emph{περίσσειας} $\mathrm{ex}(\gamma)$,
\begin{equation}\label{eq:excess}
R(y)=x+\gamma\,\frac{z}{\tfrac1n\sum_jz_j/x_j},\qquad z=(y-x)_+,
\end{equation}
η οποία διατηρεί κάθε κορυφή στο ύψος της αρχικής της μάζας ή υψηλότερα και σταθεροποιεί τον μέσο λόγο περίσσειας $\tfrac1n\sum_j(y_j-x_j)/x_j$ στο $\gamma\in\{1,4,16\}$. Οι τελεστές της Ενότητας~\ref{sec:operators} δίνουν τις δύο παραλλαγές των χειρογράφων: $S=EW\,y$ (έλξη με κλιμάκωση διαφυγής) και $S=Wy$ («κλασικό PPR» όπως διατυπώθηκε, δηλαδή έλξη χωρίς κλιμάκωση διαφυγής)· προσθέτουμε τον τελεστή ώθησης $W^{\top}$, τον ακατέργαστο πίνακα γειτνίασης $P$ και τον $EP$.

\subsection{Η αγκυρωμένη οικογένεια δυνάμεων}\label{sec:ap}

Σε μια κανονικοποιημένη κλίμακα η περίσσεια $y_i-x_i$ δεν έχει σταθερή αφετηρία: οι μισές περίπου κορυφές πέφτουν κάτω από την αρχική τους μάζα μετά την κανονικοποίηση $L_1$, και ο πολλαπλασιασμός του $y$ με μια σταθερά μεταβάλλει το $(y_i-x_i)/x_i$ μη ομοιόμορφα. Ο λόγος $q_i=y_i/x_i=1+(y_i-x_i)/x_i$ δεν έχει αυτό το πρόβλημα, διότι μια ολική αλλαγή κλίμακας πολλαπλασιάζει όλα τα $q_i$ με τον ίδιο παράγοντα. Γι' αυτό εντάσσουμε τον νόμο στην οικογένεια
\begin{equation}\label{eq:ap}
\frac{\alpha_i}{1-\alpha_i}=\exp\Big(c+\eta\Big(\log q_i-\frac1n\sum_{j}\log q_j\Big)\Big),\qquad q_i=\frac{y_i}{\tilde x_i},\qquad \tilde x=(1-\tau)\,x+\frac{\tau}{n}\one,
\end{equation}
όπου $c$ είναι η μοναδική σταθερά με $\frac1n\sum_i\alpha_i=\alpha_0$ (ο μέσος όρος του αριστερού μέλους είναι γνησίως αύξων ως προς $c$· τον επιλύουμε με προστατευμένη επανάληψη Newton). Ονομάζουμε την~\eqref{eq:ap} \emph{αγκυρωμένο αναλογικό νόμο} $\AP(\eta)$. Έχει τρεις ιδιότητες που την καθιστούν τον δικαιότερο έλεγχο της αρχής: για $\eta=0$ \emph{είναι} ο βαθμονομημένος βαθμωτός κανόνας ($\alpha_i\equiv\alpha_0$)· είναι αναλλοίωτη ως προς την ολική κλίμακα του $y$· και για $\eta=1$ θέτει τον λόγο πιθανοτήτων απόσβεσης ανάλογο του $q_i$, δηλαδή υλοποιεί την αρχή με το ολικό της επίπεδο καθορισμένο από τον κανόνα του~\cite{konstantellos26} (η προσθετική μονάδα που χωρίζει το $q_i$ από το $(y_i-x_i)/x_i$ δεν είναι αναγνωρίσιμη σε κανονικοποιημένη κλίμακα και απορροφάται από την αγκύρωση). Ο ρυθμιστής $\tau$ υλοποιεί την ομαλοποίηση του παρονομαστή που προτάθηκε στα χειρόγραφα· $\tau=0$ εκτός αν δηλώνεται διαφορετικά. Ο $\AP(\eta)$ εκτελείται είτε μέχρι το αυτοσυνεπές σταθερό σημείο του (κοινή οκνηρή επανάληψη των $y$ και $\alpha$) είτε ως εξομάλυνση με σταθερό $\alpha$.

\subsection{Εκδοχές εκτέλεσης}\label{sec:modes}

Η \emph{Εκδοχή~(α)} εκτελεί μια υλοποίηση ή τον $\AP(\eta)$ από τη θερμή εκκίνηση για $T$ βήματα ή έως τη σύγκλιση. Η \emph{Εκδοχή~(β)} υπολογίζει πρώτα ένα διάνυσμα βάσης $y^{\mathrm{base}}$ --- τον βαθμονομημένο κανόνα ή δύο επαναλήψεις της μετάδοσης μηνυμάτων του~\cite{knz14} --- το εμφυτεύει στην κλίμακα περίσσειας, $y^{(0)}=x+\gamma z$ με $z=\N(y^{\mathrm{base}})/\big(\tfrac1n\sum_j\N(y^{\mathrm{base}})_j/x_j\big)$, εκτελεί $K$ βήματα ενός νόμου σε κάποιον τελεστή ώστε να λάβει το $y^{\mathrm{prop}}$, και συνδυάζει
\begin{align}
y^{\mathrm{final}}&=(1-\beta)\,\N\big(y^{\mathrm{base}}\big)+\beta\,\N\big(y^{\mathrm{prop}}\big) &&\text{(συνδυασμός βαθμολογιών)},\label{eq:blend}\\
r^{\mathrm{final}}&=(1-\beta)\,r\big(y^{\mathrm{base}}\big)+\beta\,r\big(y^{\mathrm{prop}}\big) &&\text{(συνδυασμός τάξεων)},\nonumber
\end{align}
όπου $r(\cdot)$ συμβολίζει το διάνυσμα των τάξεων. Η ονομαστική Εκδοχή~(β) των χειρογράφων --- Υλοποίηση~Β στην κλίμακα περίσσειας, $K=1$ βήμα, $\beta=0.5$ --- αναφέρεται ως \emph{Εκδοχή~(β) ως έχει}.

\section{Δυναμική θεωρία}\label{sec:theory}

Σε όλη την ενότητα $T\ge0$ κατά στοιχείο, οι ανισότητες μεταξύ διανυσμάτων ισχύουν κατά συνιστώσα, και ο $T$ λέγεται \emph{υποστοχαστικός κατά γραμμές} αν $T\one\le\one$· οι $W$ και $EW$ είναι υποστοχαστικοί κατά γραμμές, ενώ ο τελεστής ώθησης $W^{\top}$ έχει αθροίσματα στηλών ίσα με τη μονάδα και φασματική ακτίνα ίση με τη μονάδα. Εκτός αν δηλώνεται διαφορετικά, οι υλοποιήσεις εκτελούνται στην ακατέργαστη κλίμακα με $\alpha_{\max}=1$. Το Σχήμα~\ref{fig:dynamics} και ο Πίνακας~\ref{tab:dynamics} (Παράρτημα~\ref{app:tables}) παρουσιάζουν τις αντίστοιχες μετρήσεις στο σύνολο ανάπτυξης.

\begin{figure}[!t]
\centering
\includegraphics[width=\linewidth]{fig1_dynamics.pdf}
\caption{Δυναμική των υλοποιήσεων (σύνολο ανάπτυξης). (α)~Κλάσμα κορυφών στον τετριμμένο κλάδο $y_i=x_i$ για τις Υλοποιήσεις Α, Β και Γ στον $EW$ στην ακατέργαστη κλίμακα: υπό ομοιόμορφη αρχική κατανομή η Α καταρρέει στο $t=1$ (Πρόταση~A) και η Β απορροφάται μέσα σε λίγα βήματα (Πρόταση~B). (β)~Μέσος συντελεστής $\rho$ του Spearman κατά τις επαναλήψεις, με χρώματα όπως στο~(α)· συνεχής γραμμή: $EW$, ακατέργαστη κλίμακα· διακεκομμένη: $EW$, κλίμακα περίσσειας $\gamma=16$· στικτή: Υλοποίηση~Γ στον ακατέργαστο πίνακα γειτνίασης $P$. Καμία διαμόρφωση δεν φθάνει τον βαθμονομημένο κανόνα του~\cite{konstantellos26}. (γ)~Φασματικές ακτίνες των Ιακωβιανών στο θετικό αυτοσυνεπές σημείο για $T=\lambda W$· αμφότερες υπερβαίνουν τη μονάδα για $1<\lambda<2$ (Πρόταση~C)· οι κύκλοι είναι αριθμητικές τιμές. (δ)~Η ενημέρωση της Υλοποίησης~Γ, $S\mapsto L_2(x,S)/x$: πέφτει κάτω από την αρχική μάζα όταν $S<x$ και είναι φθίνουσα για $S<(\sqrt2-1)x$ (Πρόταση~D).}
\label{fig:dynamics}
\end{figure}

\subsection{Κατάρρευση και απορρόφηση σε συσταλτικούς τελεστές}

\begin{prop}{A}[κατάρρευση της Υλοποίησης Α σε πεπερασμένο χρόνο]
Έστω $T$ υποστοχαστικός κατά γραμμές, $x=\tfrac1n\one$, και έστω ότι η Υλοποίηση~Α εκτελείται από οποιοδήποτε $y^{(0)}\ge x$. Τότε το $M_t=\max_iy^{(t)}_i$ ικανοποιεί $M_{t+1}\le\max\big(\tfrac1n,\,M_t-\tfrac1n\big)$, οπότε $y^{(t)}=x$ για κάθε $t\ge\lceil nM_0\rceil-1$. Από τη θερμή εκκίνηση $y^{(0)}=x+Tx$ η κατάρρευση συμβαίνει στο $t=1$.
\end{prop}

\begin{proof}
Κάθε επανάληψη ικανοποιεί $y^{(t)}\ge x$, διότι η ενημέρωση είναι $\max(x,S-x)$. Επειδή $T\ge0$ και $T\one\le\one$, ισχύει $S_i=(Ty^{(t)})_i\le M_t$, άρα $y^{(t+1)}_i=\max\big(\tfrac1n,S_i-\tfrac1n\big)\le\max\big(\tfrac1n,M_t-\tfrac1n\big)$. Επαγωγικά $M_t\le\max\big(\tfrac1n,M_0-\tfrac tn\big)$, που ισούται με $\tfrac1n$ μόλις $t\ge nM_0-1$· μαζί με το $y^{(t)}\ge x$ αυτό δίνει $y^{(t)}=x$. Το τετριμμένο σημείο είναι τότε σταθερό, διότι από το $Tx\le x$ έπεται $S-x\le0<x$. Για τη θερμή εκκίνηση, $y^{(0)}=x+Tx\le2x$, οπότε $M_0\le2/n$ και αρκεί ένα βήμα.
\end{proof}

Υπό ομοιόμορφη αρχική κατανομή, το \TrivAOneUnif\% των κορυφών βρίσκεται στον τετριμμένο κλάδο μετά από ένα βήμα στον $EW$ (Σχήμα~\ref{fig:dynamics}α), και σε όλα τα ζεύγη στιγμιοτύπου--τελεστή που ελέγχθηκαν η μεγαλύτερη σχετική περίσσεια που απομένει μετά από ένα βήμα είναι \CollapseExcess, δηλαδή σφάλμα στρογγυλοποίησης. Υπό ετερογενείς αρχικές κατανομές μια κορυφή παραμένει ενεργή μόνο όσο $(Ty)_i>2x_i$, πράγμα που απαιτεί γείτονες με πολύ περισσότερη αρχική μάζα από την ίδια την κορυφή· το \ActADirEnd\% των κορυφών παραμένει ενεργό υπό αρχικές κατανομές Dirichlet(1). Η συνθήκη ευνοεί κορυφές με μικρή αρχική μάζα σε σχέση με τους γείτονές τους, όχι κορυφές που είναι καλά συνδεδεμένες.

\begin{prop}{B}[απορρόφηση και τετραγωνική συστολή της Υλοποίησης Β]
Έστω $T$ υποστοχαστικός κατά γραμμές και έστω ότι εκτελείται η Υλοποίηση~Β. (i)~Αν $y_i^{(t)}\le x_i$ για κάποιο $t$, τότε $y_i^{(s)}=x_i$ για κάθε $s>t$. (ii)~Αν $x=\xi\one$ είναι ομοιόμορφη και $m_t=\max_i\big(y_i^{(t)}-\xi\big)_+$, τότε
\begin{equation}\label{eq:quad}
m_{t+1}\le\frac{m_t^2}{\xi+m_t},
\end{equation}
οπότε το $m_t$ φθίνει προς το μηδέν, τετραγωνικά μόλις $m_t\ll\xi$, και $y^{(t)}\to x$.
\end{prop}

\begin{proof}
(i) Αν $y_i^{(t)}\le x_i$, ο υλοποιημένος νόμος δίνει $\alpha_i^{(t+1)}=0$, άρα $y_i^{(t+1)}=x_i$ και, επαναλαμβάνοντας το επιχείρημα, $y_i^{(s)}=x_i$ για κάθε $s>t$. (ii) Έστω $e_i=(y_i^{(t)}-\xi)_+\le m_t$. Τότε $\alpha_i^{(t+1)}=e_i/(e_i+\xi)\le m_t/(m_t+\xi)$, επειδή η $e\mapsto e/(e+\xi)$ είναι αύξουσα, και $S_i\le\max_jy^{(t)}_j\le\xi+m_t$, επειδή ο $T$ είναι υποστοχαστικός κατά γραμμές. Αν $S_i\le\xi$, η ενημέρωση $(1-\alpha_i)\xi+\alpha_iS_i$ δεν υπερβαίνει το $\xi$. Διαφορετικά $y_i^{(t+1)}-\xi=\alpha_i^{(t+1)}(S_i-\xi)\le\frac{m_t}{m_t+\xi}\,m_t$. Αυτή είναι η~\eqref{eq:quad}. Αφού $m^2/(\xi+m)<m$ για $m>0$ και $m_t\le m_0$, η περίσσεια συστέλλεται τουλάχιστον γεωμετρικά με λόγο $m_0/(\xi+m_0)<1$, και τετραγωνικά όταν $m_t\ll\xi$· οι κορυφές κάτω από το $\xi$ επιστρέφουν στο $\xi$ στο επόμενο βήμα λόγω του~(i).
\end{proof}

Στον $EW$ υπό ομοιόμορφη αρχική κατανομή, το \TrivBFiveUnif\% των κορυφών έχει απορροφηθεί έως το $t=5$ και το \TrivBTwentyUnif\% έως το $t=20$. Η απορρόφηση είναι ακριβώς αυτό που η χρονικά υστερημένη μορφή σχεδιάστηκε να αποφύγει: τα χειρόγραφα την εισήγαγαν επειδή δεν έχει κατώφλι στο $2x_i$, όμως σε έναν συσταλτικό τελεστή η περίσσεια από την οποία τροφοδοτείται συρρικνώνεται σε κάθε βήμα.

\subsection{Αστάθεια της αυτοσυνεπούς λύσης}

\begin{prop}{C}[οι θετικές αυτοσυνεπείς λύσεις είναι απωθητικές]
Έστω $T\ge0$ και έστω $y>0$ ένα πλήρως ενεργό σταθερό σημείο, $y=Ty-x$ με $Ty>2x$.
\begin{enumerate}[label=(\roman*),leftmargin=2em]
\item $\rho(T)>1$. Ειδικότερα, τέτοιο σημείο δεν υπάρχει όταν $\rho(T)\le1$, π.χ.\ για τους $W$, $EW$ και $W^{\top}$.
\item Οι Ιακωβιανές των απεικονίσεων των Υλοποιήσεων Α και Β στο $y$ είναι $J_A=T$ και $J_B=\diag(x/y)+\diag(\one-x/y)\,T$, και $\rho(J_A)>1$, $\rho(J_B)>1$· το σημείο είναι ασταθές και για τις δύο υλοποιήσεις.
\item Για $T=\lambda W$ με $W$ στοχαστικό κατά γραμμές και $x=\xi\one$, το διάνυσμα $y=x/(\lambda-1)$ είναι πλήρως ενεργό σταθερό σημείο αν και μόνο αν $1<\lambda<2$, και τότε
\begin{equation}\label{eq:rhoJB}
\rho(J_A)=\lambda,\qquad \rho(J_B)=-\lambda^2+3\lambda-1,
\end{equation}
αμφότερες μεγαλύτερες της μονάδας στο $(1,2)$.
\end{enumerate}
\end{prop}

\begin{proof}
Χρησιμοποιούμε το φράγμα Collatz--Wielandt~\cite[Θεώρ.~8.1.26]{hornjohnson13}: αν $A\ge0$, $v>0$ και $Av\ge cv$, τότε $\rho(A)\ge c$. (i)~Η $Ty=y+x$ δίνει $(Ty)_i/y_i=1+x_i/y_i$, του οποίου το ελάχιστο υπερβαίνει τη μονάδα. (ii)~Στο ενεργό σύνολο η απεικόνιση της Α είναι η $y\mapsto Ty-x$, με Ιακωβιανή $T$. Η απεικόνιση της Β είναι $y'_i=x_i^2/y_i+(1-x_i/y_i)S_i$ με $S=Ty$· η ρητή παράγωγός της ως προς $y_i$ είναι $x_i(S_i-x_i)/y_i^2=x_i/y_i$ στο σταθερό σημείο (όπου $S_i-x_i=y_i$), και η παράγωγός της ως προς $S_i$ είναι $1-x_i/y_i$, γεγονός που δίνει την $J_B$. Και οι δύο Ιακωβιανές είναι μη αρνητικές, επειδή $y=S-x>x$. Επιπλέον $J_Ay=y+x>y$ και
\[
J_By=x+(\one-x/y)\circ(y+x)=y+x\circ\frac{y-x}{y}>y,
\]
άρα $\rho(J_A)>1$ και $\rho(J_B)>1$ από το ίδιο φράγμα. Ένα σταθερό σημείο στο οποίο η Ιακωβιανή έχει ιδιοτιμή μέτρου μεγαλύτερου της μονάδας είναι ασταθές. (iii)~Αφού $W\one=\one$, το $y=c\one$ λύνει την $y=\lambda Wy-\xi\one$ αν και μόνο αν $c=\xi/(\lambda-1)$, που είναι θετικό αν και μόνο αν $\lambda>1$ και υπερβαίνει το $\xi$ (ισοδύναμα $S>2x$) αν και μόνο αν $\lambda<2$. Τότε $x/y=(\lambda-1)\one$ και $J_B=(\lambda-1)I+(2-\lambda)\lambda W$. Για $a,b\ge0$ η φασματική ακτίνα του $aI+bW$ είναι $a+b\,\rho(W)=a+b$, που δίνει την~\eqref{eq:rhoJB}· τέλος $-\lambda^2+3\lambda-2=(\lambda-1)(2-\lambda)>0$ στο $(1,2)$.
\end{proof}

Για $\lambda={}$\LamValOne{}, \LamValTwo{} και \LamValThree{} οι αριθμητικά υπολογισμένες φασματικές ακτίνες του $J_B$ είναι \RhoJBOne, \RhoJBTwo{} και \RhoJBThree, όπως προβλέπει η~\eqref{eq:rhoJB} (Σχήμα~\ref{fig:dynamics}γ), και μια διαταραχή σχετικού μεγέθους $10^{-9}$ μεγεθύνεται μέχρι οι επαναλήψεις να εγκαταλείψουν το σταθερό σημείο. Η Πρόταση~C και η Παρατήρηση~\ref{rem:scale} εκφράζουν το ίδιο πράγμα από δύο πλευρές: η ενεργή λύση απαιτεί επεκτατικό τελεστή, και σε επεκτατικό τελεστή είναι απωθητική. Οι συσταλτικοί τελεστές επομένως καταρρέουν (Προτάσεις A και~B), ενώ στους επεκτατικούς, όπως ο $P$, οι επαναλήψεις στην ακατέργαστη κλίμακα είτε αυξάνονται κατά μήκος του κυρίαρχου ιδιοδιανύσματος του $T$ είτε πέφτουν στον τετριμμένο κλάδο.

\subsection{Η Υλοποίηση Γ είναι διαφορετικός νόμος}

\begin{prop}{D}[η Υλοποίηση Γ συστέλλεται σε σημείο κάτω από την αρχική κατανομή]
Έστω $\gamma=1$, $T$ υποστοχαστικός κατά γραμμές, $f(x_i,s)=(x_i^2+s^2)/(x_i+s)=L_2(x_i,s)$, και έστω οποιοδήποτε $y^{(0)}>0$.
\begin{enumerate}[label=(\roman*),leftmargin=2em]
\item $f(x_i,s)\ge(2\sqrt2-2)\,x_i$ για κάθε $s\ge0$, με ισότητα αν και μόνο αν $s=(\sqrt2-1)x_i$· η $f$ είναι φθίνουσα ως προς $s$ στο $[0,(\sqrt2-1)x_i]$ και αύξουσα στη συνέχεια· και $f(x_i,s)<x_i$ αν και μόνο αν $0<s<x_i$.
\item Για $X\ge\max\big(\max_ix_i,\max_iy^{(0)}_i\big)$ το υπερορθογώνιο $\mathcal B=\{y:(2\sqrt2-2)x\le y\le X\one\}$ περιέχει κάθε επανάληψη από το $t=1$ και μετά, και η απεικόνιση της Υλοποίησης~Γ είναι συστολή στη νόρμα $\infty$ επί του $\mathcal B$. Η επανάληψη επομένως συγκλίνει σε μοναδικό σταθερό σημείο $y^C$ από κάθε θετική εκκίνηση.
\item Κάθε κορυφή με $0<S_i<x_i$ στο $y^C$ έχει $y^C_i<x_i$, και μια απομονωμένη κορυφή ($S_i=0$, $y^C_i=x_i$) κατατάσσεται πάνω από κάθε μη απομονωμένη κορυφή με την ίδια αρχική μάζα και $S_i<x_i$.
\end{enumerate}
\end{prop}

\begin{proof}
(i) Με $r=s/x_i$, το $f/x_i=(1+r^2)/(1+r)$ έχει παράγωγο $(r^2+2r-1)/(1+r)^2$, η οποία μηδενίζεται στο $r=\sqrt2-1$, είναι αρνητική πριν και θετική μετά· η ελάχιστη τιμή είναι $(4-2\sqrt2)/\sqrt2=2\sqrt2-2$. Επίσης $f<x_i$ αν και μόνο αν $s^2<x_is$. (ii) Από το~(i) κάθε επανάληψη είναι τουλάχιστον $(2\sqrt2-2)x$. Ο μέσος Lehmer βρίσκεται μεταξύ των ορισμάτων του, άρα $y^{(t+1)}_i\le\max(x_i,S_i)\le\max(x_i,\max_jy^{(t)}_j)\le X$, όπου χρησιμοποιήθηκε το $S_i\le\max_jy^{(t)}_j$. Η απεικόνιση $F(y)_i=f\big(x_i,(Ty)_i\big)$ έχει Ιακωβιανή $\diag(c)\,T$ με $c_i=\partial_sf=1-2x_i^2/(x_i+S_i)^2\in(-1,1)$ για $S_i>0$. Στο $\mathcal B$, μια κορυφή με τουλάχιστον έναν γείτονα έχει $S_i\in\big[(2\sqrt2-2)(Tx)_i,\;X\big]$, ένα συμπαγές διάστημα εντός του $(0,\infty)$, οπότε $|c_i|\le\kappa_i<1$· μια απομονωμένη κορυφή έχει μηδενική γραμμή. Άρα $\|\diag(c)T\|_\infty\le\max_i\kappa_i<1$ στο κυρτό σύνολο $\mathcal B$, η $F$ είναι εκεί συστολή από την ανισότητα μέσης τιμής, και το θεώρημα σταθερού σημείου του Banach δίνει μοναδικό σταθερό σημείο και σύγκλιση. Τα υπερορθογώνια για μεγαλύτερο $X$ περιέχουν εκείνα για μικρότερο $X$, άρα το σταθερό σημείο δεν εξαρτάται από το $X$. (iii) Έπεται από το~(i) εφαρμοσμένο στο σταθερό σημείο.
\end{proof}

Τρεις τυχαίες εκκινήσεις φθάνουν στο ίδιο σταθερό σημείο με απόκλιση έως \CStartsDiff{} (Σχήμα~\ref{fig:dynamics}δ). Για γενικό $\gamma>0$ η ενημέρωση είναι $(x_i^2+\gamma S_i^2)/(x_i+\gamma S_i)$, και το σκέλος~(iii) ισχύει αυτούσιο, διότι $x_i^2+\gamma S_i^2<x_i(x_i+\gamma S_i)$ αν και μόνο αν $S_i<x_i$. Η Υλοποίηση~Γ είναι επομένως καλά ορισμένο δυναμικό σύστημα, δεν είναι όμως ο αναλογικός νόμος: ο λόγος πιθανοτήτων της, $S_i/x_i$, συγκρίνει το \emph{εισερχόμενο} σήμα με την αρχική μάζα αντί για τη συσσωρευμένη περίσσεια, και το σταθερό της σημείο αντιστρέφει τη σειρά ασθενώς συνδεδεμένων και απομονωμένων κορυφών.

\subsection{Το εμπόδιο προσήμου}

\begin{prop}{E}[εμπόδιο προσήμου]
(i) Για κάθε ενημέρωση $u_i=(1-\alpha_i)x_i+\alpha_iS_i$ ισχύει $\partial u_i/\partial\alpha_i=S_i-x_i$. Ένας νόμος που δίνει μεγαλύτερο $\alpha_i$ στις καλύτερα συνδεδεμένες κορυφές --- όπως ο $\alpha_i=1-x_i/y_i$ --- τις ανεβάζει μόνο όπου $S_i>x_i$ και τις κατεβάζει όπου $S_i<x_i$.
(ii) Έστω $T$ υποστοχαστικός κατά γραμμές, $\alpha\in[0,1)^n$, $D=\diag(\alpha)$, και έστω $y(\alpha)=(I-DT)^{-1}(I-D)x$ η γραμμική διάχυση ανά κόμβο με $S=Ty$. Τότε
\begin{equation}\label{eq:sign}
\frac{\partial y}{\partial\alpha_k}=(S_k-x_k)\,G e_k,\qquad G=(I-DT)^{-1}\ge0,\qquad g_{kk}\ge g_{jk}\ \ \text{για κάθε }j,
\end{equation}
άρα η $\partial(y_k-y_j)/\partial\alpha_k=(S_k-x_k)(g_{kk}-g_{jk})$ έχει το πρόσημο του $S_k-x_k$ ή μηδενίζεται: η αύξηση του $\alpha_k$ μπορεί να ανεβάσει την κορυφή~$k$ στην κατάταξη μόνο αν $S_k>x_k$.
(iii) Για $T=EW$ και ομοιόμορφη αρχική κατανομή, το $S=EW\,x$ έχει $S_i=\Pesc(i)/n<x_i$ για κάθε κορυφή με $\Pesc(i)<1$.
\end{prop}

\begin{proof}
Το (i) είναι άμεσο. (ii) Οι γραμμές του $DT$ αθροίζουν το πολύ σε $\max_k\alpha_k<1$, άρα $\rho(DT)<1$ και $G=\sum_{m\ge0}(DT)^m\ge0$ (σειρά Neumann). Παραγωγίζοντας την $y=(I-D)x+DTy$ ως προς $\alpha_k$ λαμβάνουμε $(I-DT)\,\partial y/\partial\alpha_k=(S_k-x_k)e_k$. Για την ιδιότητα της στήλης, έστω $v=Ge_k/g_{kk}$ (σημειώνεται ότι $g_{kk}\ge1$), οπότε $v_k=1$ και $v_j=\sum_l(DT)_{jl}v_l$ για $j\ne k$. Έστω ότι $\mu=\max_jv_j>1$ και $A=\{j:v_j=\mu\}\not\ni k$. Για $j\in A$ ισχύει $\mu=\sum_l(DT)_{jl}v_l\le\mu\sum_l(DT)_{jl}\le\mu$, άρα κάθε γραμμή του $DT$ με δείκτη στο $A$ αθροίζει στη μονάδα και έχει φορέα εντός του $A$· ο κύριος υποπίνακας $(DT)_{AA}$ είναι τότε στοχαστικός κατά γραμμές με φασματική ακτίνα ίση με τη μονάδα, σε αντίφαση με το $\rho\big((DT)_{AA}\big)\le\rho(DT)<1$. Επομένως $v\le\one$, δηλαδή $g_{jk}\le g_{kk}$. (iii) $W\one=\one$ στις μη απομονωμένες κορυφές, άρα $(EW\,\tfrac1n\one)_i=\Pesc(i)/n$.
\end{proof}

Το σκέλος~(iii) σημαίνει ότι στην ίδια την αρχική κατανομή κάθε κορυφή βρίσκεται στη λάθος πλευρά της συνθήκης προσήμου· στο σταθερό σημείο του κανόνα, σε ένα αραιό στιγμιότυπο Erdős--Rényi, το ίδιο ισχύει για το \SparseFracSbelowX\% των κορυφών. Για την κανονικοποιημένη διάχυση της Ενότητας~\ref{sec:operators}, οι πεπερασμένες διαφορές συμφωνούν με τον κανόνα προσήμου (η μεταβολή της βαθμολογίας μιας κορυφής μετά την αύξηση του $\alpha$ της έχει το πρόσημο του $S_k-x_k$) σε \SignAgree{} από 20 τυχαία επιλεγμένες κορυφές. Η Ενότητα~\ref{sec:sign} δείχνει ότι το ποσοστό των κορυφών με $S_i>x_i$ προβλέπει, μεταξύ στιγμιοτύπων, αν η βέλτιστη απόσβεση ανά κόμβο συμφωνεί με τον νόμο.

\subsection{Η παραλλαγή έλξης χωρίς κλιμάκωση διαφυγής}

\begin{prop}{F}[εκφυλισμός της διάχυσης έλξης υπό ομοιόμορφες αρχικές κατανομές]
Έστω $W$ στοχαστικός κατά γραμμές χωρίς απομονωμένες κορυφές, $x$ ομοιόμορφη και $S=Wy$. Τότε κάθε επανάληψη των Υλοποιήσεων Α, Β και Γ από τη θερμή εκκίνηση $x+Wx$, στην ακατέργαστη κλίμακα, στην $L_1$ ή στην κλίμακα περίσσειας, είναι πολλαπλάσιο του $x$: όλες οι κορυφές ισοπαλούν και $\rho=0$.
\end{prop}

\begin{proof}
$Wx=x$. Αν $y^{(t)}=cx$ τότε $S=cx$, και οι τρεις ενημερώσεις δίνουν $\max(1,c-1)\,x$, $\big((1-a)+ac\big)x$ με $a=(1-1/c)_+$, και $\frac{1+\gamma c^2}{1+\gamma c}\,x$· οι τρεις απεικονίσεις κλίμακας στέλνουν πολλαπλάσια του $x$ σε πολλαπλάσια του $x$. Η θερμή εκκίνηση είναι $2x$.
\end{proof}

Στο σύνολο ανάπτυξης, το \PullZeroPct\% των αυτόνομων εκτελέσεων με $S=Wy$ υπό ομοιόμορφη αρχική κατανομή λαμβάνει ακριβώς $\rho=0$, και ο μέσος όρος του συνόλου είναι \PullMean. Η ουσιαστική παραλλαγή «κλασικού PPR» είναι ο τελεστής ώθησης $W^{\top}$, τον οποίο χρησιμοποίησε το~\cite{konstantellos26} και τον οποίο συμπεριλαμβάνουμε παντού.

\subsection{Επανεξέταση της Πρότασης~2}\label{sec:prop2}

Τα συνοδευτικά χειρόγραφα διατυπώνουν, ως Πρόταση~2, ότι (1)~υπό ομοιόμορφη αρχική κατανομή το $\alpha_i=1-1/(ny_i)$ είναι γνησίως μονότονη και αναλλοίωτη σε μετατοπίσεις συνάρτηση μόνο της προσβασιμότητας $y_i$, και (2)~καθώς η αρχική κατανομή εντοπίζεται, $|\partial_{x_i}((y_i-x_i)/x_i)|=y_i/x_i^2\to\infty$ και $\alpha_i\to1$ για κάθε κορυφή που δεν είναι σπόρος. Το σκέλος~(1) είναι ορθό ως προς τη μονοτονία, δεν συνεπάγεται όμως πλεονέκτημα κατάταξης, διότι κατά την Πρόταση~E ένας μονότονος νόμος εξακολουθεί να δρα με το πρόσημο του $S_i-x_i$· επιπλέον η απεικόνιση δεν είναι αναλλοίωτη σε μετατοπίσεις (η πρόσθεση μιας σταθεράς στο $y$ μεταβάλλει τα $\alpha_i$ μη ομοιόμορφα). Το σκέλος~(2) είναι ορθή παράγωγος, και η συνέπειά του --- ο κορεσμός --- είναι αυτό που παρατηρούμε υπό αρχικές κατανομές σπόρων. Υπό ομαλές ετερογενείς αρχικές κατανομές η αποτυχία είναι διαφορετικού είδους, και η ακόλουθη ταυτότητα την καθιστά ακριβή.

\begin{prop}{$2'$}[ταυτότητα λογαριθμικής διασποράς]
Έστω $z_i=y_i-x_i>0$ η εξωτερική μάζα και $r_i=z_i/x_i$ ο αναλογικός λόγος. Επί των κορυφών, έστω ότι τα $\log z$ και $\log x$ έχουν τυπικές αποκλίσεις $\sigma_z,\sigma_x>0$ και συσχέτιση $\phi$, και ότι το $\log r$ δεν είναι σταθερό. Τότε
\begin{equation}\label{eq:logdisp}
\mathrm{corr}(\log r,\log z)=\frac{\sigma_z-\phi\,\sigma_x}{\sqrt{\sigma_z^2+\sigma_x^2-2\phi\,\sigma_z\sigma_x}},\qquad
\mathrm{corr}(\log r,-\log x)=\frac{\sigma_x-\phi\,\sigma_z}{\sqrt{\sigma_z^2+\sigma_x^2-2\phi\,\sigma_z\sigma_x}} .
\end{equation}
Καθώς $\sigma_x/\sigma_z\to\infty$, η πρώτη τείνει στο $-\phi$ και η δεύτερη στη μονάδα: όταν η λογαριθμική διασπορά της αρχικής κατανομής κυριαρχεί έναντι εκείνης της εξωτερικής μάζας, ο λόγος διατάσσει τις κορυφές κατά $1/x_i$ και δεν διατηρεί καμία πληροφορία για το $z$ πέρα από τη συσχέτισή του με την αρχική κατανομή.
\end{prop}

\begin{proof}
$\log r=\log z-\log x$, άρα
\[
\mathrm{Var}(\log r)=\sigma_z^2+\sigma_x^2-2\phi\sigma_z\sigma_x,\quad
\mathrm{Cov}(\log r,\log z)=\sigma_z^2-\phi\sigma_z\sigma_x,\quad
\mathrm{Cov}(\log r,-\log x)=\sigma_x^2-\phi\sigma_z\sigma_x.
\]
Η διαίρεση δίνει την~\eqref{eq:logdisp}· διαιρώντας αριθμητές και παρονομαστές με $\sigma_x$ λαμβάνουμε τα όρια.
\end{proof}

Στην πραγματική τιμή αναφοράς του συνόλου αρχικών κατανομών (Ενότητα~\ref{sec:prior}), η συσχέτιση τάξεων του $r$ με το $1/x$ αυξάνεται από \RinvDirTen{} υπό Dirichlet(10) σε \RinvDirSmall{} υπό Dirichlet(0.3) ($\log_{10}\kappa\approx{}$\KappaDirSmall{}), ενώ η συσχέτισή του με το $y^*$ μειώνεται από \RyDirTen{} σε \RyDirSmall. Υπό αρχικές κατανομές σπόρων η πλειονότητα των κορυφών που δεν είναι σπόροι μοιράζεται μία τιμή αρχικής μάζας, οπότε το $r$ παραμένει διατακτικά συνεπές με το $y^*$ (\RySeedFive{} και \RySeedOne{} για πέντε και έναν σπόρο), αλλά κορέννυται: το \SatSeedOne\% των κορυφών λαμβάνει $\alpha_i>0.99$ υπό έναν μόνο σπόρο, πράγμα που καταργεί την επανεκκίνηση για το μεγαλύτερο μέρος του γραφήματος. Η ομαλοποίηση του παρονομαστή $\tilde x=(1-\tau)x+\tau/n$ αντιμετωπίζει και τις δύο αποτυχίες, και για $\tau=1$ απομακρύνει πλήρως την αρχική κατανομή από τον παρονομαστή.

\FloatBarrier
\section{Πειραματικό πρωτόκολλο}\label{sec:protocol}

\subsection{Σύνολα στιγμιοτύπων}\label{sec:panels}

Ο Πίνακας~\ref{tab:panels} παρουσιάζει τα πέντε σύνολα στιγμιοτύπων (\NTotal{} στιγμιότυπα). Τα συνθετικά στιγμιότυπα προέρχονται από \NFamilies{} οικογένειες: Erdős--Rényi~\cite{erdos59} (μέσος βαθμός~5), Barabási--Albert~\cite{barabasi99} ($m=2$), Watts--Strogatz~\cite{watts98} ($k=6$, πιθανότητα επανασύνδεσης $0.1$), τετραγωνικά πλέγματα και ορθογώνια πλέγματα αναλογίας $1{:}4$ (τα πλέγματα στα οποία το~\cite{konstantellos26} ανέφερε το μεγαλύτερο περιθώριο βελτίωσης), στοχαστικά μοντέλα τεσσάρων μπλοκ~\cite{holland83}, γραφήματα πυρήνα--περιφέρειας (core--periphery)~\cite{borgatti00}, ομοιόμορφα τυχαία δένδρα και γραφήματα αλτήρα (barbell). Οι πιθανότητες των ακμών λαμβάνονται ομοιόμορφα από ένα από πέντε εύρη, $R_1=[0.01,0.06]$, $R_2=[0.06,0.18]$, $R_3=[0.18,0.40]$, $R_4=[0.40,0.70]$ και $R_5=[0.70,0.95]$· μόνο στα γραφήματα Erdős--Rényi το $\Pescbar$ καλύπτει έτσι το διάστημα από \ERPbarMin{} έως \ERPbarMax.

Τα σύνολα ανάπτυξης και ελέγχου διασταυρώνουν τις εννέα οικογένειες, τα πέντε εύρη και δύο αρχικές κατανομές (ομοιόμορφη και Dirichlet(1)) με δύο ανεξάρτητες επαναλήψεις η καθεμία, για $n=100$· το σύνολο κλίμακας επαναλαμβάνει τον σχεδιασμό για $n=900$ με μία επανάληψη. Το σύνολο αρχικών κατανομών κρατά σταθερά το γράφημα και τις πιθανότητες των ακμών και μεταβάλλει μόνο την αρχική κατανομή, σε δέκα τύπους: ομοιόμορφη· Dirichlet($\theta$) με $\theta\in\{10,3,1,0.3\}$· λογαριθμοκανονική με $\sigma\in\{0.5,1,2\}$· και εξομαλυμένες κατανομές σπόρων, στις οποίες 5 ή 1 σπόροι μοιράζονται το $1-\varepsilon$ της μάζας και το $\varepsilon=0.1$ κατανέμεται ομοιόμορφα. Καλύπτει πέντε οικογένειες και τα εύρη $R_2$--$R_4$. Το σύνολο πραγματικών δικτύων χρησιμοποιεί τη λέσχη καράτε του Zachary~\cite{zachary77}, το δίκτυο συνεμφανίσεων των Αθλίων (Les Misérables)~\cite{knuth93} και το δίκτυο γυναικών του Νότου του Davis~\cite{davis41}, καθώς και δύο δείγματα χιονοστιβάδας (snowball) 500 κορυφών από καθένα από τα Cora, CiteSeer, Amazon-Photo και MS-Academic-CS της συλλογής του~\cite{shchur18}. Κάθε πραγματικό γράφημα λαμβάνει τέσσερα σχήματα πιθανοτήτων --- ομοιόμορφες στα $[0.05,0.2]$, $[0.2,0.5]$ και $[0.5,0.9]$, και έναν συμμετρικοποιημένο σταθμισμένο καταρράκτη (weighted cascade) $p_{ij}=1/\max(\deg i,\deg j)$~\cite{kkt03} --- και αμφότερες τις αρχικές κατανομές.

\TabPanels

\subsection{Πραγματική τιμή αναφοράς}\label{sec:gt}

Κάθε στιγμιότυπο φέρει δύο ανεξάρτητες εκτιμήσεις Monte Carlo του~\eqref{eq:target}, τις $y^*_A$ και $y^*_B$, η καθεμία από \MGTsp{} δειγματοληπτημένους κόσμους. Οι κόσμοι δειγματοληπτούνται κατά δέσμες που στοιβάζονται σε ένα μπλοκ-διαγώνιο γράφημα, έτσι ώστε μία μόνο κλήση υπολογισμού συνεκτικών συνιστωσών να επισημαίνει όλους τους κόσμους μιας δέσμης. Σε μικρά γραφήματα ο εκτιμητής συμφωνεί με την ακριβή απαρίθμηση όλων των $2^{|E|}$ κόσμων εντός του σφάλματος δειγματοληψίας (μέγιστη τιμή $|z|$ ίση με \MCExactZ). Κάθε μέθοδος αξιολογείται ως προς το $y^*_A$· κάθε μαντείο και κάθε πείραμα «αληθούς $y^*$» χρησιμοποιεί μόνο το $y^*_B$, ώστε καμία προσαρμοσμένη ποσότητα να μη βλέπει το δείγμα αξιολόγησης. Η συσχέτιση τάξεων μεταξύ $y^*_B$ και $y^*_A$, η \emph{αναφορά διαιρεμένου δείγματος}, μετρά την αξιοπιστία της πραγματικής τιμής αναφοράς. Επειδή και τα δύο δείγματα έχουν θόρυβο, ένας ακριβής προβλεπτής του $y^*$ θα έπιανε περίπου την τετραγωνική ρίζα αυτής της αναφοράς, οπότε η αναφορά δεν αποτελεί αυστηρό ανώτατο όριο.

\subsection{Μέθοδοι}\label{sec:methods}

\emph{Αναφορές.} Ο βαθμονομημένος κανόνας (κανονικοποιημένη διάχυση στον $EW$ με $\alpha_0$)~\cite{konstantellos26}· η ίδια διάχυση με το καθιερωμένο $\alpha=0.85$· το βαθμονομημένο εξατομικευμένο PageRank ώθησης του~\cite{konstantellos26} ($W^{\top}$ με $\alpha=\clip(0.241+0.859\,\Pescbar,0.4,1)$)· η απλοϊκή τοπική απόσβεση $\alpha_i=\Pesc(i)$· και δύο επαναλήψεις της μετάδοσης μηνυμάτων των Karrer, Newman και Zdeborová~\cite{knz14}.

\emph{Μαντεία.} Το βαθμωτό μαντείο επιλέγει την καλύτερη από τις 40 τιμές $\alpha\in\{0.025,0.05,\dots,1\}$ για το $y^*_B$. Το μαντείο ανά κόμβο εκτελεί ανάβαση κατά συντεταγμένες επί του $\alpha\in[0,1]^n$ ξεκινώντας από την τιμή του βαθμωτού μαντείου, με πλέγμα 14 τιμών ανά συντεταγμένη, το πολύ τρία περάσματα και χρονικό όριο (120\,s, ή 240\,s στα \NPernodeScale{} στιγμιότυπα τετραγωνικού πλέγματος, ορθογώνιου πλέγματος, Erdős--Rényi και Barabási--Albert του συνόλου κλίμακας στα οποία εκτελείται). Επειδή προσαρμόζεται στο $y^*_B$ και αξιολογείται στο $y^*_A$, εκτιμά το περιθώριο βελτίωσης της απόσβεσης ανά κόμβο εκτός δείγματος· εντός δείγματος αποτελεί κάτω φράγμα του βελτίστου.

\emph{Ο νόμος με το αληθές $y^*$.} Η διάχυση στον $EW$ με $\alpha_i=1-x_i/y^*_{B,i}$ και ο αγκυρωμένος νόμος~\eqref{eq:ap} με $q=y^*_B/x$. Αυτά απομακρύνουν κάθε πηγή σφάλματος εκτίμησης από την αρχή.

\subsection{Επιλογή, πάγωμα και στατιστική ανάλυση}\label{sec:stats}

Όλες οι σχεδιαστικές επιλογές της αρχής αναζητήθηκαν μόνο στο σύνολο ανάπτυξης: η υλοποίηση· ο τελεστής ($EW$, $W$, $W^{\top}$, $P$, $EP$)· η κλίμακα (ακατέργαστη, $L_1$, περίσσειας με $\gamma\in\{1,4,16\}$)· ο αριθμός βημάτων· το $\eta\in\{-1,-0.5,-0.25,0,0.25,0.5,1,2\}$· η βάση (βαθμονομημένος κανόνας ή μετάδοση μηνυμάτων)· το $K\in\{1,2,5,\infty\}$· το $\beta\in\{0.25,0.5,0.75,1\}$· και το είδος του συνδυασμού. Η αναζήτηση κάλυψε \CfgSA{} αυτόνομες, \CfgAP{} αγκυρωμένες, \CfgPB{} διαμορφώσεις αγκυρωμένης εξομάλυνσης και \CfgPBX{} διαμορφώσεις εξομάλυνσης με υλοποίηση, συνολικά \CfgTotal. Οι νικήτριες διαμορφώσεις, οι ονομαστικές ρυθμίσεις των χειρογράφων και οι αντιστοιχισμένοι έλεγχοι που περιγράφονται παρακάτω παγώθηκαν ως \NFrozen{} διαμορφώσεις (Πίνακας~\ref{tab:configs}, Παράρτημα~\ref{app:tables}) σε αρχείο του οποίου η σύνοψη SHA-256 αρχίζει με \texttt{\SelHash} (\SelTime), πριν εξεταστεί οποιοδήποτε αποτέλεσμα των συνόλων ελέγχου, αρχικών κατανομών, πραγματικών δικτύων ή κλίμακας. Κάθε παγωμένη διαμόρφωση αξιολογήθηκε στη συνέχεια μία φορά σε κάθε ανεξάρτητο σύνολο.

\emph{Αντιστοιχισμένοι έλεγχοι.} Για να απομονωθεί η συνεισφορά του νόμου ανά κόμβο μέσα σε μια ροή επεξεργασίας, κάθε παγωμένη διαμόρφωση που εκτελεί μια υλοποίηση εκτελείται εκ νέου με το διάνυσμα $\alpha$ αντικατεστημένο σε κάθε βήμα από τη μέση τιμή του (το \emph{δίδυμο μέσης τιμής}) ή από μια τυχαία μετάθεσή του· ο τελεστής, η εμφύτευση, το $K$ και ο συνδυασμός παραμένουν σταθερά. Τέσσερις από τις παγωμένες διαμορφώσεις είναι τέτοιοι έλεγχοι για τον καλύτερο εξομαλυντή με υλοποίηση: το δίδυμο μέσης τιμής του, η μετάθεσή του και η ίδια ροή επεξεργασίας με σταθερό $\alpha\in\{0.5,0.9\}$.

\emph{Στατιστική ανάλυση.} Για κάθε μέθοδο αναφέρουμε τον μέσο όρο του $\rho$ στο σύνολο και τη διαφορά κατά ζεύγη ως προς τον βαθμονομημένο κανόνα, με διάστημα εμπιστοσύνης 95\% bootstrap εκατοστημορίων από 20\,000 επαναδειγματοληψίες~\cite{efron93}· τις νίκες, ισοπαλίες και ήττες (Ν/Ι/Η, όπου $|\Delta\rho|\le10^{-9}$ μετρά ως ισοπαλία)· και έναν αμφίπλευρο έλεγχο προσημασμένων τάξεων Wilcoxon~\cite{wilcoxon45} στα μη ισόπαλα ζεύγη. Σε κάθε σύνολο οι τιμές $p$ των \NFrozen{} παγωμένων διαμορφώσεων διορθώνονται κατά Holm~\cite{holm79}· οι αναφορές και τα μαντεία αναφέρονται χωρίς διόρθωση. Ο έλεγχος Wilcoxon αφορά τη θέση των διαφορών κατά ζεύγη, ενώ το διάστημα bootstrap τη μέση τιμή τους· τα δύο μπορεί να διαφωνούν όταν λίγες μεγάλες διαφορές κυριαρχούν στη μέση τιμή, και αναφέρουμε αμφότερα. Οι διαφορές τυπώνονται με τρία δεκαδικά ψηφία, και ένα προσημασμένο μηδέν όπως το $-0.000$ δηλώνει διαφορά απόλυτης τιμής κάτω από $5\times10^{-4}$. Δευτερεύουσες μετρικές είναι ο συντελεστής $\tau$-b του Kendall~\cite{kendall38}, η επικάλυψη των κορυφαίων 10\% και το NDCG@10\%~\cite{jarvelin02} με συνάφεια $y^*_A$.

\emph{Υλοποίηση.} Όλες οι μέθοδοι έχουν υλοποιηθεί σε Python με NumPy και αραιούς πίνακες SciPy~\cite{harris20,virtanen20}, και τα συνθετικά γραφήματα παράγονται με το NetworkX~\cite{hagberg08}. Κάθε τυχαία επιλογή λαμβάνει σπόρο από μια σύνοψη SHA-256 του ρόλου της, και όλοι οι χρόνοι είναι διάμεσοι χρόνοι εκτέλεσης σε μία διεργασία, στο ίδιο μηχάνημα.

\FloatBarrier
\section{Αποτελέσματα}\label{sec:results}

\subsection{\texorpdfstring{Αναπαραγωγή του~\cite{konstantellos26} και το περιθώριο βελτίωσης ανά κόμβο}{Αναπαραγωγή του κανόνα βαθμονόμησης και το περιθώριο βελτίωσης ανά κόμβο}}\label{sec:replication}

Ο Πίνακας~\ref{tab:replication} αναπαράγει τη βαθμονόμηση του~\cite{konstantellos26} στα σύνολα ανάπτυξης και ελέγχου. Υπό αρχικές κατανομές Dirichlet ο κανόνας ανακτά το \RepPctDir\% του βαθμωτού μαντείου εντός δείγματος και υπερέχει της καλύτερης σταθεράς που επιλέγεται εκ των υστέρων (\RepRuleDir{} έναντι \RepBestConstDir{} για $\alpha={}$\RepBestConstAlphaDir{}), όπως αναφέρθηκε εκεί. Υπό ομοιόμορφες αρχικές κατανομές ανακτά μόνο το \RepPctUnif\%: το μέσο $\alpha$ του μαντείου είναι \RepOracleAlphaUnif{} έναντι \RepRuleAlpha{} του κανόνα, η μία σταθερά $\alpha={}$\RepConstAlphaUnif{} επιτυγχάνει \RepConstUnif{} έναντι \RepRuleUnif{} του κανόνα, και ακόμη και το καθιερωμένο $\alpha=0.85$ επιτυγχάνει \RepFixedUnif. Η βαθμονόμηση του~\cite{konstantellos26} εξαρτάται επομένως από την αρχική κατανομή: προσαρμόστηκε εκεί όπου η βέλτιστη απόσβεση είναι υψηλή, και υπεραποσβένει όταν η αρχική κατανομή είναι ομοιόμορφη.

\TabReplication

Ο Πίνακας~\ref{tab:headroom} μετρά εκ νέου το περιθώριο βελτίωσης της απόσβεσης ανά κόμβο, με το μαντείο να προσαρμόζεται στο $y^*_B$ και να αξιολογείται στο $y^*_A$. Εκτός δείγματος, το μαντείο ανά κόμβο υπερέχει του βαθμωτού μαντείου κατά \HeadAll{} (εντός δείγματος κατά \HeadApparent, άρα το \HeadOverfitPct\% του φαινομενικού κέρδους είναι υπερπροσαρμογή στο δείγμα προσαρμογής), κατά \HeadGrid{} στα τετραγωνικά πλέγματα, \HeadRect{} στα ορθογώνια πλέγματα και \HeadTree{} στα δένδρα. Το περιθώριο συγκεντρώνεται στη γωνία των πυκνών γραφημάτων με ετερογενή αρχική κατανομή (εύρος $R_5$ με αρχικές κατανομές Dirichlet: \HeadCorner{} κατά μέσο όρο επί των οικογενειών), όπου κανένα βαθμωτό $\alpha$ δεν κατατάσσει καλά· κέρδος του μεγέθους που ανέφερε το~\cite{konstantellos26} εμφανίζεται μόνο σε μεμονωμένα πυκνά στιγμιότυπα πλεγμάτων (μέγιστο \HeadMax). Η απόσβεση ανά κόμβο έχει όντως περιθώριο βελτίωσης, αλλά μικρότερο και πιο εντοπισμένο από όσο υπέθετε το κίνητρο της αρχής.

\TabHeadroom

\subsection{Επιβεβαιωτική αξιολόγηση}\label{sec:confirm}

Ο Πίνακας~\ref{tab:confirm} παρουσιάζει κάθε παγωμένη διαμόρφωση στο ανεξάρτητο σύνολο ελέγχου μαζί με τις αναφορές και τα μαντεία, και το Σχήμα~\ref{fig:modeb}α συνοψίζει τις κύριες μεθόδους και στα τέσσερα ανεξάρτητα σύνολα.

\emph{Εκδοχή~(α).} Κάθε αυτόνομη υλοποίηση υστερεί του κανόνα: οι καλύτερες διαμορφώσεις των Α, Β και Γ μεταβάλλουν το $\rho$ κατά \TestAbestD, \TestBbestD{} και \TestCbestD{} ($p$ κατά Holm = \TestAbestHolm, \TestBbestHolm{} και \TestCbestHolm). Οι δύο επιλογές που αφορούν συγκεκριμένη αρχική κατανομή αποδίδουν μόνο υπό την κατανομή για την οποία επιλέχθηκαν. Η Υλοποίηση~Γ στον ακατέργαστο πίνακα γειτνίασης, η καλύτερη επιλογή για αρχικές κατανομές Dirichlet, επιτυγχάνει \TestCdirOnDir{} υπό Dirichlet αλλά \TestCdirOnUnif{} υπό ομοιόμορφες· η Υλοποίηση~Β στον τελεστή ώθησης, η καλύτερη επιλογή για ομοιόμορφες αρχικές κατανομές, επιτυγχάνει \TestBunifOnUnif{} υπό ομοιόμορφες αλλά \TestBunifOnDir{} υπό Dirichlet.

\emph{Ο αγκυρωμένος νόμος.} Ο $\AP(\eta=1)$ μεταβάλλει το $\rho$ κατά \TestAPoneD{} \TestAPoneCI{} ($p$ κατά Holm = \TestAPoneHolm) και ο $\AP(\eta=0.25)$ κατά \TestAPquarterD{} ($p$ κατά Holm = \TestAPquarterHolm). Το Σχήμα~\ref{fig:sign}γ δείχνει ολόκληρη την καμπύλη ως προς $\eta$. Υπό αρχικές κατανομές Dirichlet κορυφώνεται στο $\eta={}$\EtaBestDir{}, δηλαδή στον ίδιο τον βαθμωτό κανόνα, τόσο στο σύνολο ανάπτυξης όσο και στο σύνολο ελέγχου, και στο σύνολο ελέγχου φθίνει από \EtaZeroTestDir{} για $\eta=0$ σε \EtaOneTestDir{} για $\eta=1$ και \EtaTwoTestDir{} για $\eta=2$· υπό ομοιόμορφες αρχικές κατανομές είναι επίπεδη (εύρος \EtaUnifRange{} για $\eta\in[-1,2]$). Η μεταβολή ανά κόμβο προς την κατεύθυνση του νόμου δεν βοηθά ποτέ.

\emph{Εκδοχή~(β) με τον νόμο.} Η Εκδοχή~(β) ως έχει μεταβάλλει το $\rho$ κατά \TestPBXnomBD{} ($p$ κατά Holm = \TestPBXnomBHolm): ένα βήμα της Υλοποίησης~Β στην κλίμακα περίσσειας αλλοιώνει την κατάταξη βάσης, και ο συνδυασμός 50\% δεν την αποκαθιστά. Το αγκυρωμένο αντίστοιχό της ($\AP(\eta=1)$, $K=1$, $\beta=0.5$) μεταβάλλει το $\rho$ κατά \TestPBnomD{} ($p$ κατά Holm = \TestPBnomHolm) και η αυτοσυνεπής εκδοχή με συνδυασμό τάξεων κατά \TestPBoneD. Η καλύτερη αγκυρωμένη εξομάλυνση που βρέθηκε στο σύνολο ανάπτυξης χρησιμοποιεί $\eta=-0.25$, δηλαδή κινείται \emph{αντίθετα} προς τον νόμο, και δεν διακρίνεται από τον κανόνα (\TestPBbestDfour· ισοπαλία στα περισσότερα στιγμιότυπα, Ν/Ι/Η \TestPBbestWTL).

\emph{Ο νόμος με το αληθές $y^*$.} Ο ορισμός $\alpha_i=1-x_i/y^*_i$ με το ανεξάρτητο δείγμα $y^*_B$ απομακρύνει πλήρως το σφάλμα εκτίμησης. Επιτυγχάνει \TestTrueProp{} (\TestTruePropD), και ο αγκυρωμένος νόμος με το αληθές $y^*$ επιτυγχάνει \TestTrueAPone{} (\TestTrueAPoneD). Η γνώση του $y^*$ δεν καθιστά τον νόμο χρήσιμο: το εμπόδιο βρίσκεται στον νόμο και όχι στα δεδομένα εισόδου του.

\emph{Εκδοχή~(β) με δεύτερο τελεστή.} Οι παγωμένες διαμορφώσεις που βελτιώνουν πράγματι τον κανόνα συνδυάζουν όλες τη βαθμονομημένη διάχυση με έναν \emph{διαφορετικό} τελεστή διάδοσης. Η Υλοποίηση~Γ στον ακατέργαστο πίνακα γειτνίασης ως εξομαλυντής ($\gamma=4$, $K=5$, $\beta=0.5$) κερδίζει \TestPBXbestD{} \TestPBXbestCI{} ($p$ κατά Holm = \TestPBXbestHolm). Οι αντιστοιχισμένοι έλεγχοί της κερδίζουν το ίδιο ή περισσότερο: το δίδυμο μέσης τιμής \TestCmeanD, ένα σταθερό $\alpha=0.5$ στον $P$ \TestLinHalfD{} και ένα σταθερό $\alpha=0.9$ στον $P$ με συνδυασμό τάξεων \TestLinNineD{} \TestLinNineCI{} ($p$ κατά Holm = \TestLinNineHolm, Ν/Ι/Η \TestLinNineWTL). Ο τελευταίος φθάνει $\bar\rho={}$\TestLinNine{}, δηλαδή το \TestLinNinePctOracle\% του βαθμωτού μαντείου· καλύπτει το \TestLinNineGapScalar\% του χάσματος μεταξύ κανόνα και βαθμωτού μαντείου και το \TestLinNineGapNode\% του χάσματος έως το μαντείο ανά κόμβο. Η εξομάλυνση της βάσης μετάδοσης μηνυμάτων δύο επαναλήψεων με τον βαθμωτό κανόνα ($\eta=0$) κερδίζει \TestMPzeroD, ενώ ο νόμος στην ίδια ροή επεξεργασίας ($\eta=1$) κερδίζει λιγότερο (\TestMPoneD). Ο έλεγχος μετάθεσης κερδίζει \TestCpermD{} κατά μέσο όρο αλλά χάνει στα περισσότερα στιγμιότυπα (Ν/Ι/Η \TestCpermWTL· $p$ κατά Holm = \TestCpermHolm), παράδειγμα διαφωνίας μεταξύ μέσης τιμής και θέσης κατά Wilcoxon.

\TabConfirm

\begin{figure}[!t]
\centering
\includegraphics[width=\linewidth]{fig2_sign_obstruction.pdf}
\caption{Το εμπόδιο προσήμου στα δεδομένα. (α)~Συσχέτιση τάξεων μεταξύ της απόσβεσης $\alpha^*$ του μαντείου ανά κόμβο και του νόμου $1-x_i/y^*_i$, έναντι του κλάσματος των κορυφών των οποίων το σήμα γειτόνων υπερβαίνει την αρχική τους μάζα στο σταθερό σημείο του κανόνα (ένα σημείο ανά στιγμιότυπο ανάπτυξης ή ελέγχου· μαύρο: μέσοι όροι ανά κλάση). (β)~Διαφορά στο $\rho$ μεταξύ του νόμου με το αληθές $y^*$ και του βαθμονομημένου κανόνα. (γ)~Μέσο $\rho$ του αγκυρωμένου νόμου $\AP(\eta)$ στα σύνολα ανάπτυξης (διακεκομμένη) και ελέγχου (συνεχής)· το $\eta=0$ (συνεχής κατακόρυφη γραμμή) είναι ο βαθμωτός κανόνας και το $\eta=1$ (στικτή γραμμή) ο αναλογικός νόμος.}
\label{fig:sign}
\end{figure}

\subsection{Το εμπόδιο προσήμου στα δεδομένα}\label{sec:sign}

Η Πρόταση~E προβλέπει ότι ο νόμος μπορεί να συμφωνεί με μια καλή απόσβεση ανά κόμβο μόνο σε κορυφές με $S_i>x_i$. Το Σχήμα~\ref{fig:sign}α το ελέγχει αυτό στα \SignN{} στιγμιότυπα ανάπτυξης και ελέγχου. Το κλάσμα των κορυφών με $S_i>x_i$ στο σταθερό σημείο του κανόνα προβλέπει τη συσχέτιση τάξεων μεταξύ του $\alpha^*$ του μαντείου ανά κόμβο και του νόμου (Spearman \SignSpearman, $p={}$\SignSpearmanP{}). Στα \SignLowN{} στιγμιότυπα όπου λιγότερο από το 10\% των κορυφών ικανοποιεί $S_i>x_i$ η συσχέτιση είναι \SignLowCorr{} --- η βέλτιστη απόσβεση κινείται \emph{αντίθετα} προς τον νόμο --- και ο νόμος με το αληθές $y^*$ χάνει \SignLowDrop{} στο $\rho$, δηλαδή αντιστρέφει την κατάταξη (Σχήμα~\ref{fig:sign}β). Όπου οι περισσότερες κορυφές ικανοποιούν τη συνθήκη προσήμου η συσχέτιση γίνεται θετική, αλλά δεν υπερβαίνει ποτέ το \SignMaxCorr.

\subsection{Ομοιομορφία της αρχικής κατανομής}\label{sec:prior}

Ο Πίνακας~\ref{tab:prior} και το Σχήμα~\ref{fig:prior} παρουσιάζουν το σύνολο αρχικών κατανομών, στο οποίο μεταβάλλεται μόνο η αρχική κατανομή. Η διαφορά μεταξύ $\AP(\eta=1)$ και κανόνα είναι \PriorAPUnif{} υπό ομοιόμορφη αρχική κατανομή και μειώνεται με τη διασπορά, μονότονα εντός κάθε οικογένειας κατανομών, σε \PriorAPDirTen{} υπό Dirichlet(10), \PriorAPDirSmall{} υπό Dirichlet(0.3), \PriorAPLognTwo{} υπό λογαριθμοκανονική(2) και \PriorAPSeedOne{} υπό έναν μόνο σπόρο· η συσχέτισή της κατά Spearman με το $\log_{10}\kappa(x)$ είναι \PriorSpearman{} ($p={}$\PriorSpearmanP{}). Η κατεύθυνση της Πρότασης~2 των χειρογράφων επιβεβαιώνεται, αλλά δεν υπάρχει καθεστώς πλεονεκτήματος: το «καλύτερη υπό ομοιόμορφες αρχικές κατανομές» σημαίνει «αβλαβής υπό ομοιόμορφες αρχικές κατανομές». Το Σχήμα~\ref{fig:prior}β δείχνει τον μηχανισμό που εντόπισε η Πρόταση~$2'$, και το Σχήμα~\ref{fig:prior}γ την επίδραση της ομαλοποίησης: το $\tau=0.5$ μειώνει την απώλεια (π.χ.\ \PriorTauHalfLognTwo{} υπό λογαριθμοκανονική(2)) και το $\tau=1$ την εξαλείφει (μεταξύ \PriorTauOneMin{} και \PriorTauOneMax{} σε όλους τους τύπους αρχικής κατανομής). Για $\tau=1$ ο παρονομαστής δεν εξαρτάται πλέον από το $x$, και ο νόμος ανάγεται σε μονότονο μετασχηματισμό του $y$ αγκυρωμένο στον κανόνα. Οι διαμορφώσεις διπλού τελεστή κερδίζουν υπό κάθε τύπο αρχικής κατανομής (σταθερό $\alpha=0.9$ στον $P$: \PriorLinMin{} έως \PriorLinMax· Υλοποίηση~Γ στον $P$: \PriorPBXMin{} έως \PriorPBXMax).

\TabPrior

\begin{figure}[!t]
\centering
\includegraphics[width=\linewidth]{fig3_prior_uniformity.pdf}
\caption{Ομοιομορφία της αρχικής κατανομής (σύνολο αρχικών κατανομών). (α)~Διαφορά κατά ζεύγη $\AP(\eta=1)$ μείον κανόνας έναντι του δυναμικού εύρους $\log_{10}\kappa(x)$ της αρχικής κατανομής, ένα σημείο ανά στιγμιότυπο. (β)~Μέση συσχέτιση τάξεων του λόγου $r_i=(y^*_i-x_i)/x_i$ με το $1/x_i$ και με το $y^*_i$ ανά τύπο αρχικής κατανομής, κατά σειρά αυξανόμενης διασποράς (D: Dirichlet, LN: λογαριθμοκανονική). (γ)~Επίδραση της ομαλοποίησης $\tilde x=(1-\tau)x+\tau/n$ στη διαφορά $\AP(\eta=1)$ μείον κανόνας για ομοιόμορφες, ομαλές ετερογενείς και σπόρων αρχικές κατανομές.}
\label{fig:prior}
\end{figure}

\subsection{Απόδοση: ο νόμος έναντι του διδύμου μέσης τιμής του}\label{sec:attribution}

Ο Πίνακας~\ref{tab:attribution} και το Σχήμα~\ref{fig:modeb}β συγκρίνουν κάθε παγωμένη διαμόρφωση που εκτελεί μια υλοποίηση με το δίδυμο μέσης τιμής της και με τη μετάθεσή της. Στο σύνολο ελέγχου ο καλύτερος εξομαλυντής είναι ελαφρώς \emph{χειρότερος} από το δίδυμό του (\AttTestPBXD{} \AttTestPBXCI· $p$ κατά Holm = \AttTestPBXHolm), και για $n=900$ οι δύο είναι ίσοι (\AttScalePBXD{} \AttScalePBXCI). Στο σύνολο πραγματικών δικτύων ο νόμος διατηρεί μια μικρή θετική επίδραση μέσα στον καλύτερο εξομαλυντή (\AttRealPBXD{} \AttRealPBXCI· $p$ κατά Holm = \AttRealPBXHolm) και μέσα στον εξομαλυντή για ομοιόμορφες αρχικές κατανομές (\AttRealPBXunifD· $p$ κατά Holm = \AttRealPBXunifHolm)· αυτές είναι οι μόνες θετικές επιδράσεις του νόμου που επιβιώνουν της διόρθωσης, και σε αμφότερες τις περιπτώσεις η ροή επεξεργασίας με σταθερό $\alpha=0.9$ είναι ακόμη καλύτερη (Πίνακας~\ref{tab:transfer}). Οι αυτόνομοι νόμοι υστερούν των διδύμων τους, συχνά σημαντικά: η Υλοποίηση~Γ στον ακατέργαστο πίνακα γειτνίασης χάνει \AttTestCdirD{} στο σύνολο ελέγχου, \AttRealCdirD{} στα πραγματικά δίκτυα και \AttScaleCdirD{} για $n=900$, επειδή η αντικατάσταση της απόσβεσης ανά κόμβο από τη μέση τιμή τη μετατρέπει σε απλή αποσβεννύμενη διάδοση στον $P$. Οι έλεγχοι μετάθεσης βρίσκονται κάτω από τα δίδυμα μέσης τιμής σε κάθε γραμμή: η αντιστοίχιση τιμών απόσβεσης σε κορυφές έχει σημασία, αλλά η αντιστοίχιση του νόμου υπερέχει μιας ομοιόμορφης μόνο στις δύο μικρές επιδράσεις των πραγματικών δικτύων που αναφέρθηκαν. Η μεταβολή ανά κόμβο που εισάγει ο νόμος είναι, στην καλύτερη περίπτωση, ουδέτερη.

\TabAttribution

\begin{figure}[!t]
\centering
\includegraphics[width=\linewidth]{fig4_modeb_and_controls.pdf}
\caption{Η Εκδοχή~(β) και οι έλεγχοί της. (α)~Μέσο $\rho$ των αναφορών, των διαμορφώσεων που χρησιμοποιούν τον αναλογικό νόμο και των αντιστοιχισμένων ελέγχων στα τέσσερα ανεξάρτητα σύνολα (στικτή γραμμή: βαθμονομημένος κανόνας). (β)~Διαφορά κατά ζεύγη μεταξύ κάθε παγωμένης διαμόρφωσης που εκτελεί μια υλοποίηση και του διδύμου μέσης τιμής της (ίδια ροή επεξεργασίας, με το $\alpha$ αντικατεστημένο από τη μέση τιμή του ανά βήμα), με διαστήματα bootstrap 95\%· τα τρίγωνα σημειώνουν διαφορές κάτω από $-0.10$, των οποίων οι τιμές αναγράφονται.}
\label{fig:modeb}
\end{figure}

\subsection{Μεταφορά σε πραγματικά δίκτυα και μεγαλύτερα γραφήματα}\label{sec:transfer}

Ο Πίνακας~\ref{tab:transfer} παρουσιάζει τις παγωμένες διαμορφώσεις στα σύνολα πραγματικών δικτύων και κλίμακας, κανένα από τα οποία δεν έπαιξε ρόλο στην επιλογή. Στα πραγματικά δίκτυα ο κανόνας επιτυγχάνει \RealRule, το βαθμωτό μαντείο \RealOracleScalar{} και το μαντείο ανά κόμβο \RealOraclePernode· το περιθώριο βελτίωσης ανά κόμβο έναντι του βαθμωτού μαντείου είναι μεγαλύτερο από ό,τι στα συνθετικά γραφήματα (\RealHeadroom· \RealHeadCora{} στα δείγματα του Cora). Ο νόμος δεν το αξιοποιεί: ο $\AP(\eta=1)$ μεταβάλλει το $\rho$ κατά \RealAPoneD{} και ο νόμος με το αληθές $y^*$ κατά \RealTruePropD. Οι διαμορφώσεις διπλού τελεστή μεταφέρονται: η Υλοποίηση~Γ στον $P$ κερδίζει \RealPBXbestD, ο έλεγχος σταθερού $\alpha$ \RealLinNineD{} \RealLinNineCI{} και η βάση μετάδοσης μηνυμάτων με εξομάλυνση από τον βαθμωτό κανόνα \RealMPzeroD. Για $n=900$ (κανόνας \ScaleRule) η εικόνα είναι η ίδια: ο έλεγχος σταθερού $\alpha$ κερδίζει \ScaleLinNineD, η Υλοποίηση~Γ στον $P$ \ScalePBXbestD, ο αγκυρωμένος νόμος μεταβάλλει το $\rho$ κατά \ScaleAPoneD{} και η καλύτερη αυτόνομη υλοποίηση κατά \ScaleCbestD. Στα \NPernodeScale{} στιγμιότυπα όπου το μαντείο ανά κόμβο εκτελείται σε αυτό το μέγεθος, υπερέχει του βαθμωτού μαντείου κατά \ScaleHeadroom· εκεί το χρονικό όριο του μαντείου είναι δεσμευτικό, οπότε το περιθώριο αυτό υποεκτιμάται.

\TabTransfer

\subsection{Δευτερεύουσες μετρικές}\label{sec:secondary}

Τα συμπεράσματα δεν εξαρτώνται από τη μετρική (Πίνακας~\ref{tab:secondary}). Στο σύνολο ελέγχου ο $\tau$ του Kendall αυξάνεται από \SecTestRuleTau{} για τον κανόνα σε \SecTestLinNineTau{} για τον έλεγχο σταθερού $\alpha$ (βαθμωτό μαντείο \SecTestOracleScalarTau), και η επικάλυψη των κορυφαίων 10\% από \SecTestRuleTop{} σε \SecTestPBXbestTop{} για τον εξομαλυντή της Υλοποίησης~Γ και σε \SecTestCmeanTop{} για το δίδυμο μέσης τιμής του. Στα πραγματικά δίκτυα οι ίδιες διαμορφώσεις αυξάνουν το NDCG@10\% από \SecRealRuleNdcg{} σε \SecRealPBXbestNdcg{} και \SecRealCmeanNdcg. Ο αγκυρωμένος νόμος και η εξομάλυνσή του υστερούν του κανόνα και στις τρεις μετρικές.

\TabSecondary

\subsection{Γιατί λειτουργεί ο συνδυασμός διπλού τελεστή}\label{sec:complement}

Αν δεν είναι ο αναλογικός νόμος αυτό που βελτιώνει την Εκδοχή~(β), τότε πρέπει να είναι ο δεύτερος τελεστής. Ο Πίνακας~\ref{tab:complement} τον απομονώνει με τη διαμόρφωση της Ενότητας~\ref{sec:explore}: πέντε βήματα της $y\leftarrow0.1\,x+0.9\,Py$ από την αρχή στον ακατέργαστο πίνακα πιθανοτήτων γειτνίασης, χωρίς καμία βαθμονομημένη βάση, επιτυγχάνουν ήδη \RawPTest{} στο σύνολο ελέγχου (\RawPTestD{} έναντι του κανόνα), \RawPReal{} στα πραγματικά δίκτυα και \RawPScale{} για $n=900$. Ο συνδυασμός τάξεων των δύο βελτιώνει και τα δύο μέλη (\BlendTest, \BlendReal{} και \BlendScale), επειδή τα σφάλματα τάξης τους είναι μόνο ασθενώς συσχετισμένα: η συσχέτιση των σφαλμάτων τάξης ανά κορυφή είναι \ErrTest{} στο σύνολο ελέγχου, \ErrReal{} στα πραγματικά δίκτυα και \ErrScale{} για $n=900$. Οι δύο τελεστές σφάλλουν με διαφορετικό τρόπο. Η διάχυση έλξης με κλιμάκωση διαφυγής σταθμίζει τις ακμές με $-\ln(1-p_{ij})$ και τις κανονικοποιεί ανά κορυφή, ενώ ο $P$ διαδίδει ακατέργαστες πιθανότητες και επιτρέπει στις κορυφές υψηλού βαθμού να συσσωρεύουν μάζα, οπότε ο συνδυασμός συμπεριφέρεται ως συνδυασμός δύο μελών με ασυσχέτιστα σφάλματα~\cite{krogh95}. Η Υλοποίηση~Γ στον $P$ λειτουργεί για τον ίδιο λόγο: η απόσβεσή της $S_i/(x_i+S_i)$ είναι κοντά στη μονάδα για τις περισσότερες κορυφές μόλις η βάση εμφυτευτεί στην κλίμακα περίσσειας, άρα είναι σχεδόν γραμμική διάδοση στον $P$, πράγμα που δείχνει ακριβώς το δίδυμο μέσης τιμής της.

\TabComplement

\subsection{Κόστος και σύγκριση με το Monte Carlo}\label{sec:cost}

Ο Πίνακας~\ref{tab:cost} και το Σχήμα~\ref{fig:cost} συγκρίνουν την ακρίβεια με τον χρόνο εκτέλεσης από άκρο σε άκρο σε μία διεργασία (συμπεριλαμβάνονται η κατασκευή τελεστών, η διάχυση, η εμφύτευση, τα βήματα και ο συνδυασμός). Για $n=100$ ο βαθμονομημένος κανόνας απαιτεί \CostTestRuleMs\,ms και ο συνδυασμός σταθερού $\alpha$ \CostTestEnsMs\,ms, ενώ μια εκτίμηση Monte Carlo 200 κόσμων κατά δέσμες απαιτεί \CostTestMCTwoHundredMs\,ms και κατατάσσει καλύτερα (\CostTestMCTwoHundredRho{} έναντι \TestLinNine)· οι 1000 κόσμοι φθάνουν \CostTestMCThousandRho{} σε \CostTestMCThousandMs\,ms. Σε αυτό το μέγεθος η δειγματοληψία είναι ταυτόχρονα φθηνότερη και ακριβέστερη από κάθε ντετερμινιστική μέθοδο της εργασίας. Για $n=900$ η σειρά του κόστους αντιστρέφεται: ο συνδυασμός απαιτεί \CostScaleEnsMs\,ms έναντι \CostScaleMCTwoHundredMs\,ms για 200 κόσμους (\CostScaleSpeedup$\times$ ταχύτερος) με ίση ακρίβεια (\ScaleLinNine{} έναντι \CostScaleMCTwoHundredRho).

Για να εντοπίσουμε το σημείο διασταύρωσης, χρονομετρήσαμε τον συνδυασμό και το Monte Carlo σε \SweepN{} στιγμιότυπα Erdős--Rényi, Barabási--Albert, Watts--Strogatz και πλέγματος με $n$ από \SweepNmin{} έως \SweepNmax{} (εύρος $R_3$, αμφότερες οι αρχικές κατανομές· Πίνακας~\ref{tab:sweep} και Σχήμα~\ref{fig:cost}γ--δ). Το μέγεθος από το οποίο ο συνδυασμός γίνεται ταχύτερος από τους 200 κόσμους εξαρτάται από την οικογένεια: $n={}$\CrossER{} στα γραφήματα Erdős--Rényi (το Monte Carlo είναι ταχύτερο για $n\le{}$\CrossERBelow{}), $n={}$\CrossBA{} στα Barabási--Albert, $n={}$\CrossGrid{} στα πλέγματα, και μεταξύ $n={}$\CrossWSBelow{} και $n={}$\CrossWS{} στα Watts--Strogatz. Δικαιότερη σύγκριση προκύπτει αν δοθεί στο Monte Carlo ο χρονικός προϋπολογισμός του συνδυασμού (ο αριθμός κόσμων που χωρούν στον ίδιο χρόνο, με παρεμβολή σε $\log k$). Η διάμεση διαφορά είναι \EqSmall{} για $n\le150$ (ο συνδυασμός υπερέχει στο \EqWinSmall\% των στιγμιοτύπων), \EqMid{} για $n=400$--$600$ και \EqLarge{} για $n\approx2000$ (υπερέχει στο \EqWinLarge\%). Η οικογένεια μετρά όσο και το μέγεθος. Η διάμεση διαφορά είναι \EqBudgetER{} στα γραφήματα Erdős--Rényi, \EqBudgetBA{} στα Barabási--Albert και \EqBudgetGrid{} στα πλέγματα, αλλά \EqBudgetWS{} στα Watts--Strogatz, όπου κάθε ντετερμινιστική μέθοδος είναι πολύ λιγότερο ακριβής από τη δειγματοληψία (διάμεσο $\rho$ \SweepWSEnsRho{} για τον συνδυασμό έναντι \SweepWSMCRho{} για 200 κόσμους). Στα πλέγματα, σε αυτό το εύρος, ο βαθμονομημένος κανόνας μόνος του είναι η ακριβέστερη ντετερμινιστική μέθοδος (\SweepGridRuleRho{} έναντι \SweepGridEnsRho).

Τα παραπάνω διορθώνουν το επιχείρημα κόστους του~\cite{konstantellos26}, το οποίο συνέκρινε τους ευρετικούς του αλγορίθμους με δειγματολήπτη περίπου μία τάξη μεγέθους βραδύτερο από έναν δειγματολήπτη κατά δέσμες. Με δειγματοληψία κατά δέσμες, οι ντετερμινιστικές διαχύσεις αξίζουν τη μεροληψία τους μόνο σε μεγαλύτερα γραφήματα και μακριά από έντονα ομαδοποιημένες, σχεδόν κρίσιμες δομές· στο σύνολο πραγματικών δικτύων ακόμη και 50 κόσμοι επιτυγχάνουν \RealMCFiftyRho{} (200 κόσμοι: \RealMCTwoHundredRho), πάνω από κάθε ντετερμινιστική μέθοδο που αξιολογήσαμε.

\begin{figure}[!t]
\centering
\includegraphics[width=\linewidth]{fig5_cost_and_monte_carlo.pdf}
\caption{Ακρίβεια, κόστος και Monte Carlo. (α,\,β)~Μέσο $\rho$ έναντι του χρόνου εκτέλεσης από άκρο σε άκρο σε μία διεργασία για $n=100$ και $n=900$ (λογαριθμική κλίμακα), με Monte Carlo κατά δέσμες για $k=50$, $200$ και $1000$ κόσμους. (γ)~Χρόνος εκτέλεσης έναντι του $n$ για τον συνδυασμό διπλού τελεστή και για το Monte Carlo κατά δέσμες (σάρωση μεγέθους, διάμεσοι). (δ)~Διαφορά στο $\rho$ μεταξύ του συνδυασμού και του Monte Carlo με τον χρονικό προϋπολογισμό του συνδυασμού, ανά οικογένεια (διάμεσοι).}
\label{fig:cost}
\end{figure}

\TabCost

\TabSweep

\subsection{Διερευνητικά: διάδοση μόνο στον ακατέργαστο πίνακα γειτνίασης}\label{sec:explore}

\emph{Η υποενότητα αυτή είναι διερευνητική και δεν αποτελούσε μέρος του παγωμένου πρωτοκόλλου.} Αφού έγιναν γνωστά τα αποτελέσματα των συνόλων ανάπτυξης και ελέγχου, εξετάσαμε αν ο δεύτερος τελεστής μπορεί να χρησιμοποιηθεί χωρίς καμία βαθμονομημένη βάση. Μεταβάλαμε τα $a\in\{0.5,0.9\}$, $K\in\{2,5,10\}$ και το βάρος συνδυασμού $\beta\in\{0.5,1\}$ για την $y\leftarrow(1-a)\,x+a\,Py$ με εκκίνηση από το $x$ --- μια περικομμένη βαθμολογία τύπου Katz~\cite{katz53} με βάρη πιθανοτήτων --- είτε μόνη της είτε σε συνδυασμό τάξεων με τον κανόνα. Η διαμόρφωση $a=0.9$, $K=5$, $\beta=0.5$ ήταν η καλύτερη και στα δύο σύνολα. Αφού επιλέχθηκε μετά την εξέτασή τους, μόνο τα σύνολα πραγματικών δικτύων και κλίμακας μπορούν να λειτουργήσουν ως ανεξάρτητος έλεγχός της: εκεί επιτυγχάνει \BlendReal{} και \BlendScale{} (Πίνακας~\ref{tab:complement}), σε συμφωνία με τον παγωμένο έλεγχο σταθερού $\alpha$ (\RealLinNine{} και \ScaleLinNine). Η διάδοση μόνο στον ακατέργαστο $P$ (\RawPReal{} και \RawPScale) είναι καλύτερη από τον κανόνα σε αυτά τα σύνολα αλλά χειρότερη από οποιονδήποτε από τους δύο συνδυασμούς. Δύο φθηνές, δομικά διαφορετικές διαδόσεις είναι καλύτερες μαζί παρά χωριστά· το αν ο βαθμονομημένος κανόνας είναι ο καλύτερος εταίρος για τον $P$ είναι ερώτημα για μια μελλοντική προκαταχωρισμένη μελέτη.

\FloatBarrier
\section{Συζήτηση}\label{sec:discussion}

\paragraph{Γιατί αποτυγχάνει η αρχή.} Τέσσερις ανεξάρτητοι μηχανισμοί, ο καθένας αποδεδειγμένος στην Ενότητα~\ref{sec:theory} και ορατός στα δεδομένα, εξηγούν το αρνητικό αποτέλεσμα. \emph{Κλίμακα:} το συζευγμένο σύστημα ισχύει στον στόχο μόνο αν $Ty^*=y^*+x$, πράγμα που δεν επιτρέπει κανένας τελεστής με φασματική ακτίνα το πολύ ίση με τη μονάδα (Παρατήρηση~\ref{rem:scale}, Πρόταση~C)· οι υλοποιήσεις επομένως επιχειρούν να φθάσουν σε μια λύση που οι κανονικοποιημένοι τελεστές δεν διαθέτουν, και είτε καταρρέουν (Προτάσεις A και~B) είτε συγκλίνουν σε διαφορετικό νόμο (Πρόταση~D). \emph{Ευστάθεια:} όπου υπάρχει η ενεργή λύση, είναι απωθητική (Πρόταση~C), οπότε δεν μπορεί να βρεθεί με επανάληψη. \emph{Πρόσημο:} ο νόμος αυξάνει την απόσβεση των καλά συνδεδεμένων κορυφών, πράγμα που αυξάνει τις βαθμολογίες τους μόνο όπου $S_i>x_i$ (Πρόταση~E)· σε κανονικοποιημένους τελεστές και αραιά καθεστώτα αυτό αποτυγχάνει για τις περισσότερες κορυφές, και η βέλτιστη απόσβεση ανά κόμβο κινείται αντίθετα προς τον νόμο (Ενότητα~\ref{sec:sign}). \emph{Αρχική κατανομή:} υπό ομαλές ετερογενείς αρχικές κατανομές ο αναλογικός λόγος μετρά το $1/x_i$ (Πρόταση~$2'$), και υπό κατανομές σπόρων κορέννυται. Η αγκύρωση και η ομαλοποίηση μπορούν να καταστήσουν τον νόμο αβλαβή --- για $\eta=0$ ή $\tau=1$, όπου δεν είναι πλέον ο νόμος --- αλλά όχι χρήσιμο. Επειδή αποτυγχάνει και ο νόμος με το \emph{αληθές} $y^*$, καμία από αυτές τις αποτυχίες δεν οφείλεται σε σφάλμα εκτίμησης.

\paragraph{Τι θα χρειαζόταν η απόσβεση ανά κόμβο.} Η απόσβεση ανά κόμβο έχει όντως περιθώριο βελτίωσης εκτός δείγματος (\HeadAll{} στα συνθετικά γραφήματα, \RealHeadroom{} στα πραγματικά δίκτυα), αλλά η απόσβεση του μαντείου είναι στην καλύτερη περίπτωση ασθενώς ευθυγραμμισμένη με τον νόμο, ακόμη και όπου ισχύει η συνθήκη προσήμου (συσχέτιση το πολύ \SignMaxCorr). Η Πρόταση~E υποδεικνύει τι πρέπει να σέβεται ένας χρήσιμος νόμος: στο σημείο λειτουργίας, η αύξηση του $\alpha_k$ ωφελεί την κορυφή~$k$ ακριβώς όταν $S_k>x_k$, οπότε κάθε κανόνας ανά κόμβο πρέπει να εξαρτάται από το πρόσημο του $S_k-x_k$ και όχι μόνο από το πόση μάζα έχει συσσωρεύσει η κορυφή. Ο σχεδιασμός και η προκαταχώριση ενός τέτοιου κανόνα αφήνονται για μελλοντική εργασία.

\paragraph{Τι προσφέρει η Εκδοχή~(β).} Η αρχιτεκτονική της Εκδοχής~(β) --- μια βαθμονομημένη διάχυση βάσης, λίγα βήματα μιας δεύτερης διάδοσης και ένας συνδυασμός --- είναι το τμήμα της πρότασης που λειτουργεί, υπό την προϋπόθεση ότι η δεύτερη διάδοση χρησιμοποιεί \emph{διαφορετικό τελεστή}. Μια σταθερή απόσβεση στον ακατέργαστο πίνακα πιθανοτήτων γειτνίασης είναι το καλύτερο δεύτερο στάδιο μεταξύ των παγωμένων διαμορφώσεων, κοστίζει $K$ γινόμενα αραιού πίνακα επί διάνυσμα, μεταφέρεται σε πραγματικά δίκτυα και στο $n=900$, και οφείλει το κέρδος της στην ασθενή συσχέτιση των σφαλμάτων τάξης των δύο τελεστών (Ενότητα~\ref{sec:complement}). Ο αναλογικός νόμος μπορεί να τοποθετηθεί σε αυτό το στάδιο χωρίς βλάβη (Υλοποίηση~Γ), αλλά το δίδυμο μέσης τιμής του είναι εξίσου καλό, και μια σταθερά είναι καλύτερη.

\paragraph{Τοποθέτηση έναντι της δειγματοληψίας.} Οι ντετερμινιστικοί ευρετικοί αλγόριθμοι κατάταξης για στοχαστικά γραφήματα συχνά αιτιολογούνται με το κόστος της δειγματοληψίας Monte Carlo. Με τη μπλοκ-διαγώνια ομαδοποίηση κατά δέσμες το κόστος αυτό είναι μικρό: σε γραφήματα εκατό κορυφών μια εκτίμηση 200 κόσμων είναι φθηνότερη και ακριβέστερη από κάθε ευρετικό αλγόριθμο της εργασίας, και ο ντετερμινιστικός συνδυασμός αξίζει τη μεροληψία του μόνο πέρα από ένα μέγεθος που εξαρτάται από την οικογένεια και μακριά από ομαδοποιημένες, σχεδόν κρίσιμες δομές όπως τα γραφήματα Watts--Strogatz. Προτείνουμε οι αξιολογήσεις τέτοιων ευρετικών να περιλαμβάνουν έναν δειγματολήπτη κατά δέσμες με τον ίδιο χρόνο εκτέλεσης.

\paragraph{Διορθώσεις του~\cite{konstantellos26}.} Δύο ευρήματα αναθεωρούν την προηγούμενη μελέτη μας. Πρώτον, η βαθμονόμησή της εξαρτάται από την αρχική κατανομή: αναπαράγεται υπό αρχικές κατανομές Dirichlet (\RepPctDir\% του βαθμωτού μαντείου) αλλά ανακτά μόνο το \RepPctUnif\% υπό ομοιόμορφη αρχική κατανομή, όπου μια σταθερή τιμή την υπερβαίνει. Δεύτερον, η σύγκριση κόστους χρησιμοποίησε δειγματολήπτη πολύ βραδύτερο από έναν δειγματολήπτη κατά δέσμες· με μέτρηση από άκρο σε άκρο, το ντετερμινιστικό πλεονέκτημα εμφανίζεται μόνο σε μεγαλύτερα γραφήματα. Το περιθώριο βελτίωσης που αποτέλεσε το κίνητρο της αρχής (\PrevHeadAll{} συνολικά, \PrevHeadGrid{} στα πλέγματα) μετράται εκ νέου σε \HeadAll{} και \HeadGrid{} εκτός δείγματος.

\section{Απειλές για την εγκυρότητα}\label{sec:limitations}

\emph{Θόρυβος της πραγματικής τιμής αναφοράς.} Με \MGTsp{} κόσμους ανά δείγμα, τα πυκνά στιγμιότυπα έχουν αναξιόπιστες τάξεις αναφοράς· στα πυκνά στιγμιότυπα Watts--Strogatz του συνόλου ανάπτυξης η αναφορά διαιρεμένου δείγματος είναι μόλις \WSDenseCeiling. Όλες οι συγκρίσεις γίνονται κατά ζεύγη, οπότε ο θόρυβος αυτός προσθέτει διακύμανση αλλά όχι μεροληψία. \emph{Ισχύς του μαντείου.} Το μαντείο ανά κόμβο είναι ανάβαση κατά συντεταγμένες σε πλέγμα 14 τιμών με το πολύ τρία περάσματα και χρονικό όριο· εντός δείγματος είναι κάτω φράγμα, και για $n=900$ στα πλέγματα το όριο είναι δεσμευτικό, οπότε το περιθώριο εκεί υποεκτιμάται, όπως και στο~\cite{konstantellos26}. \emph{Επιλογή.} Οι αναζητήσεις στο σύνολο ανάπτυξης είναι εκτεταμένες, και γι' αυτό οι νικήτριες διαμορφώσεις παγώθηκαν πριν από οποιαδήποτε ανεξάρτητη αξιολόγηση και οι ανεξάρτητοι έλεγχοι διορθώνονται κατά Holm· η διερευνητική Ενότητα~\ref{sec:explore} δεν καλύπτεται από αυτή την προστασία. \emph{Εύρος των τελεστών.} Ο νόμος ελέγχθηκε σε πέντε τελεστές και τέσσερις συμβάσεις κλίμακας, αυτόνομα και ως εξομαλυντής, με και χωρίς αγκύρωση, και με το αληθές $y^*$. Ένας διαφορετικός τελεστής θα μπορούσε να συμπεριφερθεί διαφορετικά, αλλά η Παρατήρηση~\ref{rem:scale} και η Πρόταση~C περιορίζουν κάθε τελεστή: ο στόχος ικανοποιεί το σύστημα μόνο σε επεκτατικό τελεστή, στον οποίο η λύση είναι απωθητική. \emph{Επανυλοποίηση.} Ο βαθμονομημένος κανόνας χρησιμοποιεί τις δημοσιευμένες σταθερές του~\cite{konstantellos26}· λεπτομέρειες υλοποίησης, όπως ο αριθμός επαναλήψεων και οι γεννήτριες αρχικών κατανομών, ενδέχεται να διαφέρουν από την αρχική μελέτη, όμως η αναπαραγωγή επιβεβαιώνει τα αποτελέσματά της υπό αρχικές κατανομές Dirichlet. \emph{Χρονομέτρηση.} Όλοι οι χρόνοι αφορούν Python και SciPy σε μία διεργασία στο ίδιο μηχάνημα· οι λόγοι μεταφέρονται καλύτερα από τους απόλυτους χρόνους, και μεταγλωττισμένες υλοποιήσεις οποιασδήποτε πλευράς θα μπορούσαν να μετατοπίσουν τα σημεία διασταύρωσης. \emph{Κάλυψη.} Τα σύνολα καλύπτουν εννέα συνθετικές οικογένειες, πέντε εύρη πιθανοτήτων, δέκα τύπους αρχικής κατανομής και επτά πραγματικά δίκτυα με $n\le{}$\SweepNmax{}· μεγαλύτερα γραφήματα και κατευθυνόμενα ή συσχετισμένα μοντέλα ακμών δεν έχουν ελεγχθεί.

\section{Συμπεράσματα}\label{sec:conclusion}

Η Αρχή της Αναλογικής Σχετικής Επιρροής συζευγνύει ένα βήμα διάχυσης με έναν νόμο απόσβεσης ανά κορυφή, και η άλγεβρά της είναι ακριβής: το συζευγμένο σύστημα παραγοντοποιείται ως $(y_i-x_i)(y_i+x_i-S_i)=0$. Η δυναμική και η στατιστική της συμπεριφορά είναι ωστόσο δυσμενείς. Στους κανονικοποιημένους τελεστές που χρησιμοποιούνται για κατάταξη οι υλοποιήσεις καταρρέουν στην αρχική κατανομή, η αυτοσυνεπής λύση υπάρχει μόνο σε επεκτατικούς τελεστές και εκεί είναι απωθητική, ο νόμος δρα με το πρόσημο του $S_i-x_i$ και όχι με την προσβασιμότητα της κορυφής, και οι ετερογενείς αρχικές κατανομές μετατρέπουν τον λόγο σε μέτρο του $1/x_i$. Μια προκαταχωρισμένη αξιολόγηση σε \NTotal{} στιγμιότυπα επιβεβαιώνει καθέναν από αυτούς τους μηχανισμούς: καμία διαμόρφωση που βασίζεται στον νόμο δεν υπερέχει του βαθμονομημένου βαθμωτού κανόνα του~\cite{konstantellos26}, ούτε καν με το αληθές $y^*$. Η υβριδική αρχιτεκτονική που προτάθηκε μαζί με τον νόμο είναι παρ' όλα αυτά χρήσιμη. Ο συνδυασμός της βαθμονομημένης διάχυσης με λίγα βήματα διάδοσης στον ακατέργαστο πίνακα πιθανοτήτων γειτνίασης καλύπτει το μεγαλύτερο μέρος του χάσματος έως το βαθμωτό μαντείο με κόστος λίγων αραιών γινομένων, και το κέρδος του προέρχεται από δύο τελεστές που σφάλλουν διαφορετικά. Τέλος, κάθε ισχυρισμός πρακτικού πλεονεκτήματος τέτοιων ευρετικών πρέπει να μετράται έναντι ενός δειγματολήπτη Monte Carlo κατά δέσμες, ο οποίος είναι ταχύτερος και ακριβέστερος σε μικρά γραφήματα και ξεπερνιέται μόνο πέρα από ένα μέγεθος που εξαρτάται από την οικογένεια.

\section*{Δήλωση αναπαραγωγιμότητας}

Το αρχείο που συνοδεύει την εργασία περιέχει τη βιβλιοθήκη, τη γεννήτρια των συνόλων στιγμιοτύπων, τα προγράμματα εκτέλεσης των πειραμάτων, τα προγράμματα ανάλυσης και ελέγχων, τα αποτελέσματα ανά στιγμιότυπο, το αρχείο των παγωμένων διαμορφώσεων με τη σύνοψη SHA-256 του, καθώς και τα προγράμματα που παράγουν κάθε αριθμό, πίνακα και σχήμα της εργασίας (\path{make_paper_assets.py} και \path{paper_figures.py}). Κάθε τυχαία επιλογή λαμβάνει ντετερμινιστικό σπόρο, οπότε η επανεκτέλεση ενός σταδίου αναπαράγει τα αποθηκευμένα αποτελέσματα· η ανακατασκευή όλων των μακροεντολών, πινάκων και σχημάτων από τα αποθηκευμένα αποτελέσματα απαιτεί μία εντολή. Τα αρχεία των πραγματικών δικτύων είναι εκείνα που διανέμονται με το~\cite{shchur18}. Η συνοδευτική τεκμηρίωση περιλαμβάνει την αρχιτεκτονική του κώδικα, πλήρεις αποδείξεις, περιγραφή του κώδικα ανά άρθρωμα, την πολυπλοκότητα κάθε συνιστώσας και αναλυτικό οδηγό αναπαραγωγής.

\begin{thebibliography}{99}\footnotesize
\setlength{\itemsep}{1pt}
\bibitem{kkt03} D.~Kempe, J.~Kleinberg, É.~Tardos. Maximizing the spread of influence through a social network. In \emph{Proc.\ 9th ACM SIGKDD}, 137--146, 2003.
\bibitem{hubba08} C.-Y. Lin, C.-H. Chin, H.-H. Wu, S.-H. Chen, C.-W. Ho, M.-T. Ko. Hubba: hub objects analyzer --- a framework of interactome hubs identification for network biology. \emph{Nucleic Acids Research}, 36:W438--W443, 2008.
\bibitem{cytohubba14} C.-H. Chin, S.-H. Chen, H.-H. Wu, C.-W. Ho, M.-T. Ko, C.-Y. Lin. cytoHubba: identifying hub objects and sub-networks from complex interactome. \emph{BMC Systems Biology}, 8(Suppl~4):S11, 2014.
\bibitem{konstantellos26} I.~Konstantellos. Damping from the escape probability: a free calibration rule for diffusion-based reachability ranking in uncertain graphs. Computer Engineering and Informatics Department, University of Patras, 2026. \url{https://zenodo.org/records/23211440}.
\bibitem{knz14} B.~Karrer, M.~E.~J. Newman, L.~Zdeborová. Percolation on sparse networks. \emph{Physical Review Letters}, 113:208702, 2014.
\bibitem{provanball83} J.~S. Provan, M.~O. Ball. The complexity of counting cuts and of computing the probability that a graph is connected. \emph{SIAM Journal on Computing}, 12(4):777--788, 1983.
\bibitem{potamias10} M.~Potamias, F.~Bonchi, A.~Gionis, G.~Kollios. $k$-nearest neighbors in uncertain graphs. \emph{Proceedings of the VLDB Endowment}, 3(1):997--1008, 2010.
\bibitem{pagerank99} L.~Page, S.~Brin, R.~Motwani, T.~Winograd. The PageRank citation ranking: bringing order to the Web. Technical Report 1999-66, Stanford InfoLab, 1999.
\bibitem{jeh03} G.~Jeh, J.~Widom. Scaling personalized web search. In \emph{Proc.\ 12th International World Wide Web Conference}, 271--279, 2003.
\bibitem{haveliwala03} T.~H. Haveliwala. Topic-sensitive PageRank: a context-sensitive ranking algorithm for web search. \emph{IEEE Transactions on Knowledge and Data Engineering}, 15(4):784--796, 2003.
\bibitem{gleich15} D.~F. Gleich. PageRank beyond the Web. \emph{SIAM Review}, 57(3):321--363, 2015.
\bibitem{borgs14} C.~Borgs, M.~Brautbar, J.~Chayes, B.~Lucier. Maximizing social influence in nearly optimal time. In \emph{Proc.\ 25th ACM-SIAM Symposium on Discrete Algorithms}, 946--957, 2014.
\bibitem{tim14} Y.~Tang, X.~Xiao, Y.~Shi. Influence maximization: near-optimal time complexity meets practical efficiency. In \emph{Proc.\ ACM SIGMOD}, 75--86, 2014.
\bibitem{newman23} M.~E.~J. Newman. Message passing methods on complex networks. \emph{Proceedings of the Royal Society A}, 479:20220774, 2023.
\bibitem{boldi05} P.~Boldi, M.~Santini, S.~Vigna. PageRank as a function of the damping factor. In \emph{Proc.\ 14th International World Wide Web Conference}, 557--566, 2005.
\bibitem{bressan10} M.~Bressan, E.~Peserico. Choose the damping, choose the ranking? \emph{Journal of Discrete Algorithms}, 8(2):199--213, 2010.
\bibitem{constantine10} P.~G. Constantine, D.~F. Gleich. Random alpha PageRank. \emph{Internet Mathematics}, 6(2):189--236, 2010.
\bibitem{avrachenkov18} K.~Avrachenkov, A.~Kadavankandy, N.~Litvak. Mean field analysis of personalized PageRank with implications for local graph clustering. \emph{Journal of Statistical Physics}, 173:895--916, 2018.
\bibitem{avrachenkov14} K.~Avrachenkov, R.~van der Hofstad, M.~Sokol. Personalized PageRank with node-dependent restart. In \emph{Algorithms and Models for the Web Graph (WAW 2014)}, Lecture Notes in Computer Science 8882, 23--33, Springer, 2014.
\bibitem{wu12} X.-M. Wu, Z.~Li, A.~M.-C. So, J.~Wright, S.-F. Chang. Learning with partially absorbing random walks. In \emph{Advances in Neural Information Processing Systems 25}, 2012.
\bibitem{jung19} J.~Jung, W.~Jin, H.~Park, U.~Kang. Supervised and extended restart in random walks for ranking and link prediction in networks. \emph{PLOS ONE}, 14(3):e0213857, 2019.
\bibitem{li19gpr} P.~Li, E.~Chien, O.~Milenkovic. Optimizing generalized PageRank methods for seed-expansion community detection. In \emph{Advances in Neural Information Processing Systems 32}, 2019.
\bibitem{gprgnn21} E.~Chien, J.~Peng, P.~Li, O.~Milenkovic. Adaptive universal generalized PageRank graph neural network. In \emph{Proc.\ 9th International Conference on Learning Representations}, 2021.
\bibitem{berberidis19} D.~Berberidis, A.~N. Nikolakopoulos, G.~B. Giannakis. Adaptive diffusions for scalable learning over graphs. \emph{IEEE Transactions on Signal Processing}, 67(5):1307--1321, 2019.
\bibitem{katz53} L.~Katz. A new status index derived from sociometric analysis. \emph{Psychometrika}, 18(1):39--43, 1953.
\bibitem{dwork01} C.~Dwork, R.~Kumar, M.~Naor, D.~Sivakumar. Rank aggregation methods for the Web. In \emph{Proc.\ 10th International World Wide Web Conference}, 613--622, 2001.
\bibitem{krogh95} A.~Krogh, J.~Vedelsby. Neural network ensembles, cross validation, and active learning. In \emph{Advances in Neural Information Processing Systems 7}, 231--238, 1995.
\bibitem{hornjohnson13} R.~A. Horn, C.~R. Johnson. \emph{Matrix Analysis}, 2nd edition. Cambridge University Press, 2013.
\bibitem{lehmer71} D.~H. Lehmer. On the compounding of certain means. \emph{Journal of Mathematical Analysis and Applications}, 36(1):183--200, 1971.
\bibitem{erdos59} P.~Erdős, A.~Rényi. On random graphs~I. \emph{Publicationes Mathematicae Debrecen}, 6:290--297, 1959.
\bibitem{barabasi99} A.-L. Barabási, R.~Albert. Emergence of scaling in random networks. \emph{Science}, 286(5439):509--512, 1999.
\bibitem{watts98} D.~J. Watts, S.~H. Strogatz. Collective dynamics of `small-world' networks. \emph{Nature}, 393:440--442, 1998.
\bibitem{holland83} P.~W. Holland, K.~B. Laskey, S.~Leinhardt. Stochastic blockmodels: first steps. \emph{Social Networks}, 5(2):109--137, 1983.
\bibitem{borgatti00} S.~P. Borgatti, M.~G. Everett. Models of core/periphery structures. \emph{Social Networks}, 21(4):375--395, 2000.
\bibitem{zachary77} W.~W. Zachary. An information flow model for conflict and fission in small groups. \emph{Journal of Anthropological Research}, 33(4):452--473, 1977.
\bibitem{knuth93} D.~E. Knuth. \emph{The Stanford GraphBase: A Platform for Combinatorial Computing}. Addison-Wesley, 1993.
\bibitem{davis41} A.~Davis, B.~B. Gardner, M.~R. Gardner. \emph{Deep South: A Social Anthropological Study of Caste and Class}. University of Chicago Press, 1941.
\bibitem{shchur18} O.~Shchur, M.~Mumme, A.~Bojchevski, S.~Günnemann. Pitfalls of graph neural network evaluation. Relational Representation Learning Workshop, NeurIPS 2018. arXiv:1811.05868.
\bibitem{efron93} B.~Efron, R.~J. Tibshirani. \emph{An Introduction to the Bootstrap}. Chapman \& Hall, 1993.
\bibitem{wilcoxon45} F.~Wilcoxon. Individual comparisons by ranking methods. \emph{Biometrics Bulletin}, 1(6):80--83, 1945.
\bibitem{holm79} S.~Holm. A simple sequentially rejective multiple test procedure. \emph{Scandinavian Journal of Statistics}, 6(2):65--70, 1979.
\bibitem{kendall38} M.~G. Kendall. A new measure of rank correlation. \emph{Biometrika}, 30(1--2):81--93, 1938.
\bibitem{jarvelin02} K.~Järvelin, J.~Kekäläinen. Cumulated gain-based evaluation of IR techniques. \emph{ACM Transactions on Information Systems}, 20(4):422--446, 2002.
\bibitem{harris20} C.~R. Harris et al. Array programming with NumPy. \emph{Nature}, 585:357--362, 2020.
\bibitem{virtanen20} P.~Virtanen et al. SciPy 1.0: fundamental algorithms for scientific computing in Python. \emph{Nature Methods}, 17:261--272, 2020.
\bibitem{hagberg08} A.~A. Hagberg, D.~A. Schult, P.~J. Swart. Exploring network structure, dynamics, and function using NetworkX. In \emph{Proc.\ 7th Python in Science Conference}, 11--15, 2008.
\end{thebibliography}

\appendix
\section{Πρόσθετοι πίνακες}\label{app:tables}

Ο Πίνακας~\ref{tab:configs} παρουσιάζει τις παγωμένες διαμορφώσεις με τον συμβολισμό του κώδικα, και ο Πίνακας~\ref{tab:dynamics} τα κλάσματα κλάδων στα οποία βασίζεται το Σχήμα~\ref{fig:dynamics}α.

\TabConfigs

\TabDynamics

\end{document}

